\documentclass[11pt,letterpaper]{article}

\usepackage[utf8]{inputenc}
\usepackage[T1]{fontenc}
\usepackage[margin=1in]{geometry}
\usepackage{tikz}
\usepackage{xcolor}
\usepackage{amssymb}
\usepackage{tabularx}
\usepackage{booktabs}
\usepackage{longtable}
\usepackage{enumitem}
\usepackage{hyperref}
\usepackage{titlesec}
\usepackage{fancyhdr}
\usepackage{parskip}

\definecolor{lux-red}{HTML}{8B1A1A}
\definecolor{lux-grey}{HTML}{4A4A4A}

\hypersetup{
  colorlinks=true,
  linkcolor=lux-red,
  citecolor=lux-grey,
  urlcolor=lux-red,
  pdftitle={Lux Industries Inc.\ --- IP Enforcement Memorandum},
  pdfauthor={Tridib Nandy --- Chief Technology Officer, Lux Industries Inc.},
  pdfsubject={Confidential board memorandum},
}

\titleformat{\section}{\normalfont\Large\bfseries\color{lux-grey}}{\thesection}{0.6em}{}
\titleformat{\subsection}{\normalfont\large\bfseries\color{lux-grey}}{\thesubsection}{0.6em}{}

\pagestyle{fancy}
\fancyhf{}
\fancyhead[L]{\small\color{lux-grey}\textsc{Confidential --- Board Materials --- Attorney-Client Privileged}}
\fancyhead[R]{\small\color{lux-grey}Lux Industries Inc.}
\fancyfoot[C]{\small\color{lux-grey}\thepage}
\renewcommand{\headrulewidth}{0.2pt}

\begin{document}

% ─────────────────────────────────────────────────────────────────
% Cover page
% ─────────────────────────────────────────────────────────────────
\pagestyle{empty}
\begin{tikzpicture}[remember picture,overlay]
  \fill[black] (current page.south west) rectangle (current page.north east);

  \node[anchor=north west] at ([xshift=1.4cm,yshift=-1.4cm]current page.north west) {%
    \begin{tikzpicture}[scale=0.5]
      \fill[white] (-0.7,0.55) -- (0.7,0.55) -- (0,-0.55) -- cycle;
    \end{tikzpicture}%
  };

  \node[anchor=center,text=white,font=\fontsize{30}{38}\selectfont\bfseries,text width=14cm,align=center]
    at ([yshift=3.2cm]current page.center) {LUX INDUSTRIES};

  \node[anchor=center,text=white!70,font=\fontsize{14}{18}\selectfont,text width=14cm,align=center]
    at ([yshift=1.6cm]current page.center)
    {IP Enforcement Memorandum};

  \draw[white!30,line width=0.5pt]
    ([yshift=0.5cm,xshift=-5cm]current page.center) --
    ([yshift=0.5cm,xshift=5cm]current page.center);

  \node[anchor=center,text=white!40,font=\fontsize{9}{11}\selectfont]
    at ([yshift=-0.8cm]current page.center)
    {Prepared for: Cyrus Pahlavi, President \& CEO};

  \node[anchor=center,text=white!40,font=\fontsize{9}{11}\selectfont]
    at ([yshift=-1.4cm]current page.center)
    {Hunter Dupont, Managing Partner --- Strategic Transactions};

  \node[anchor=center,text=white!40,font=\fontsize{9}{11}\selectfont]
    at ([yshift=-2.0cm]current page.center)
    {Alan Donohue, General Counsel};

  \node[anchor=center,text=white!40,font=\fontsize{9}{11}\selectfont]
    at ([yshift=-2.8cm]current page.center)
    {Prepared by: Tridib Nandy --- CTO, Lux Industries Inc.};

  \node[anchor=center,text=white!30,font=\fontsize{9}{11}\selectfont]
    at ([yshift=-3.4cm]current page.center)
    {2026-05-25};

  \node[anchor=south west,text=white!20,font=\fontsize{7}{9}\selectfont]
    at ([xshift=1.4cm,yshift=0.8cm]current page.south west)
    {Lux Industries Inc.};

  \node[anchor=south east,text=white!20,font=\fontsize{7}{9}\selectfont]
    at ([xshift=-1.4cm,yshift=0.8cm]current page.south east)
    {CONFIDENTIAL --- ATTORNEY-CLIENT PRIVILEGED};
\end{tikzpicture}
\clearpage
\pagestyle{fancy}
\setcounter{page}{1}

% ─────────────────────────────────────────────────────────────────
\begin{flushleft}
\textbf{\textsc{Memorandum}}\\[0.5em]
\begin{tabularx}{\textwidth}{lX}
\textbf{TO:}        & Cyrus Pahlavi, President \& Chief Executive Officer, Lux Industries Inc. \\
                    & Hunter Dupont, Managing Partner --- Strategic Transactions (deal lead) \\
                    & Alan Donohue, General Counsel (legal lead) \\
\textbf{CC:}        & Board of Directors, Lux Industries Inc. \\
\textbf{FROM:}      & Tridib Nandy --- Chief Technology Officer, Lux Industries Inc.; concurrently Vice President of Engineering, the Counterparty \\
\textbf{DATE:}      & 25 May 2026 \\
\textbf{RE:}        & Misappropriation of Lux foundational technology by the Counterparty (Satschel, Inc.\ d/b/a Liquidity.io, and operating affiliate Liquidity.io LLC); legal options, damages analysis, and proposed cross-license and settlement framework \\
\textbf{PRIVILEGE:} & Confidential. Attorney-client privileged. Attorney work product. Prepared in anticipation of litigation and as a board strategy document. \\
\end{tabularx}
\end{flushleft}

\vspace{0.3em}
\begin{small}
\textit{Author's disclosure of interest:} I serve concurrently as
Chief Technology Officer of Lux Industries Inc.\ and Vice President
of Engineering at the Counterparty entity referenced herein. This
memorandum is written from the Lux chair, in service of the Lux
board's decision-making. My Counterparty role gives direct
operational visibility into the matters described; I have set them
out as observed, with the understanding that any Counterparty-side
response will be developed by its own counsel and that Lux's outside
litigation counsel will need to develop independent record support
for any allegations herein before filing.
\end{small}

\vspace{0.5em}\hrule height 0.6pt\vspace{1em}

% ─────────────────────────────────────────────────────────────────
\section*{I.\ Executive Summary}

\textbf{The Counterparty has commercially deployed material portions
of Lux Industries Inc.'s patent-protected and source-available
proprietary technology since at least 2014, in production, on three
live blockchain networks, and in support of a registered ATS/BD/TA
business that is currently raising approximately \$75 million at a
\$1.1 billion valuation, with \textbf{Forum Markets} (a publicly-traded
company, formerly Ethzilla; ticker \textbf{FRMM} (NASDAQ; CIK 0001690080; principal office Palm Beach, FL; SIC 6199)) continuing as lead investor from the prior round
and verbal commitments for at least approximately half the round
secured to date --- in each case without executing the commercial
license that the governing Lux Ecosystem License v1.2
(``LEL'') and the Lux Research License with Patent Reservation v1.0
(``LRL--PR'') require.}

\textbf{Lux's direct investment in the misappropriated technology
base exceeds \$20 million in development cost recorded since 2014,
plus a multi-year contribution from senior engineering personnel
including the Lux Chief Cryptographer and the named author of this
memorandum.}

\textbf{The most striking evidentiary fact in this memorandum is the
90-day commit-ratio analysis at \S VI:} the Lux Founder and Chief
Cryptographer contributed \textbf{5,991 commits} and approximately
\textbf{93 million lines of code} (additions + deletions) across six
venture trees --- the Counterparty's primary codebase plus five
related ventures (VCC, MLC, EquityTable, OnyxPlus, and the
legacy/archive codebase including the simplici-* repositories) ---
during the most recent 90-day window, and \textbf{zero commits to
Lux's upstream repositories} in the same window. The Counterparty
reached its \$1.1B current-round valuation during that same 90 days.
The implied value-creation rate of the Lux Founder's diverted
bandwidth is approximately \textbf{\$12.2M per day}.

\textbf{Discovery posture.} This matter was identified and elevated to
the Lux Industries Inc.\ board by Mr.\ Tridib Nandy (Lux CTO; concurrent
VPE at the Counterparty), whose dual-role visibility --- \texttt{luxfi/*}
GitHub-org admin on the Lux side combined with operational seat on the
Counterparty's engineering organization --- positioned him to observe
the unauthorized integration of Lux modules into the Counterparty's
commercial stack and the absence of a corresponding executed license.
Mr.\ Nandy remains in his VPE role; his team's continued involvement is
required for the Counterparty's long-term operational success, and the
desired outcome of this matter is an \textbf{equitable settlement
preserving Counterparty operational continuity while recognising Lux's
contribution}, rather than scorched-earth litigation.

\textbf{This memorandum sets out the legal record, identifies twelve
causes of action available to Lux, quantifies damages at small /
medium / big scenarios (\$80M to \$590M for IP claims; \$180M to
\$890M with opportunity-cost component added), projects forward
enterprise value at user-scale milestones up to 100M users (Lux's
equity stake value at the rebased settlement target ranges into the
\$B+ at scale), analyzes the ``wait vs.\ settle'' time-value question,
presents a settlement triangulation matrix the deal lead (Mr.\ Dupont)
can use to negotiate, and locks the mandatory terms the President \&
CEO (Mr.\ Pahlavi) requires in any executed agreement.}

\textbf{Posture: partnership, not litigation.} The objective of this
memorandum is a \textbf{partnership agreement that regularises the
underlying IP relationship and creates a long-term commercial
alliance} between Lux Industries Inc.\ and the Counterparty.
Litigation is the alternative path available if no partnership can
be reached; the twelve available causes of action (\S V) and
litigation-cash damage ranges (\S VI, \$80M--\$890M) are documented
as \textbf{leverage and as the alternative posture}, not as the
operative ask. The corporate-lawyer framing assumes both sides
benefit from a deal that regularises the IP, preserves Counterparty
operational continuity, and aligns Lux as a substantive partner in
the Counterparty's forward enterprise value.

\textbf{Headline partnership range (equity-share).} Lux's
partnership-target range is \textbf{10\% to 50\% equity in the
Counterparty}, structured as a negotiated package of equity + cash
+ reciprocal commercial commitments. The 10\% floor reflects the
minimum substantive partnership Lux will accept; 25\% is the standard
landing zone; 35\% broadens to a deeper commercial alliance; 50\% is
the full equal-partner recapitalisation. \textbf{Higher equity tiers
carry escalating governance rights} (board observer $\to$ board seat
$\to$ governance parity) and broader reciprocal-licence packages.

\textbf{Cash and equity are inversely related to a constant
total-value target.} Target total nominal value = \$550M (= 50\% of
the \$1.1B valuation = the max-equity tier's all-equity satisfaction
of Lux's claim). Cash fills the gap at lower equity tiers:
\textbf{cash $=$ (50\% $-$ agreed equity \%) $\times$ \$1.1B}. So
50\% equity $\Rightarrow$ \$0 cash; 35\% $\Rightarrow$ \$165M cash;
25\% $\Rightarrow$ \$275M cash; 10\% $\Rightarrow$ \$440M cash. Every
tier delivers the same \$550M total nominal value to Lux. The
Counterparty chooses where in the range to land based on its
preference between existing-shareholder dilution and cash outflow.

\textbf{Substantiation at the \$1.1B valuation.} The current \$75M
raise being subscribed at \$1.1B valuation (Forum Markets, lead) is
the natural pricing comp: equity issued to Lux at the same valuation
is priced on the same per-share basis as the new-money investors,
which is the clean precedent transaction a corporate lawyer would
anchor on. See \S X for the full tier matrix with governance +
reciprocal-commitment packages and the multi-year cash sequencing
that makes the lower-equity tiers practical for the Counterparty.

\textbf{Mandatory term: Mr.\ Zach Kelling stays as Counterparty CTO
(key-person retention).} Mr.\ Kelling --- formerly Lux Founder /
Architect / Chief Cryptographer, now Counterparty CTO --- is a
prospective named defendant under the available Counts (\S V) but
is the lead key-person retention target in any partnership outcome.
Lux's equity stake is rendered substantially worthless without his
continued technical leadership; the partnership therefore pairs
mutual releases of Mr.\ Kelling individually with a multi-year
key-person retention commitment from the Counterparty. See \S XI
for the full mandatory-terms list.

\textbf{This is not a single-proposal document.} It presents data,
damage ranges, partnership-upside ranges, and negotiating floors.
The deal team triangulates to a specific outcome. Mr.\ Pahlavi
retains final approval authority on any executed agreement.

% ─────────────────────────────────────────────────────────────────
\section*{II.\ Parties and Dramatis Personae}

\subsection*{Lux Industries Inc.\ --- Decision Chain}

\begin{itemize}[leftmargin=1.5em,itemsep=0.2em]
\item \textbf{Cyrus Pahlavi} --- President \& Chief Executive Officer.
Final approval authority on any settlement, license, or litigation
decision.
\item \textbf{Hunter Dupont} --- Largest investor in Lux Industries
Inc.; Managing Partner, Strategic Transactions. Leads deal discussions
and commercial negotiations with the Counterparty.
\item \textbf{Alan Donohue} --- General Counsel. Leads legal strategy
and definitive-document drafting; provides legal support to Mr.\
Dupont throughout the negotiation.
\item \textbf{Tridib Nandy} --- Chief Technology Officer of Lux Industries
Inc.; \texttt{luxfi/*} GitHub-org admin. Author of this memorandum.
Holds concurrent role as Vice President of Engineering at the
Counterparty (\textbf{role currently active}; his team is required
for the Counterparty's long-term operational success). It was through
this concurrent visibility that Mr.\ Nandy identified the unauthorized
integration of Lux modules into the Counterparty's commercial stack
and elevated the matter to the Lux Industries Inc.\ board, occasioning
this memorandum.
\end{itemize}

\textit{Mr.\ Zach Kelling --- formerly Founder, Architect, and Chief
Cryptographer of Lux Industries Inc., and formerly a \texttt{luxfi/*}
GitHub-org admin --- has \textbf{departed Lux Industries Inc.\ and is
currently Chief Technology Officer of the Counterparty}
(Liquidity.io). He is therefore a \textbf{Counterparty officer for
purposes of this memorandum} and is named and discussed in the
Counterparty Officer Slate (below) and in \S IV.E. As a former Lux
Founder/officer and current Counterparty officer who took the
Lux-origin IP rebuild across, Mr.\ Kelling is a \textbf{prospective
named defendant} in the available causes of action (\S V), and is
the lead key-person retention target in any partnership outcome
(\S XI).}

\subsection*{Counterparty --- Officer Slate}

\begin{itemize}[leftmargin=1.5em,itemsep=0.2em]
\item \textbf{Zach Kelling} --- \textbf{Chief Technology Officer of
the Counterparty (Liquidity.io)}. Formerly Founder, Architect, and
Chief Cryptographer of Lux Industries Inc.; formerly held
\texttt{luxfi/*} GitHub-org admin. \textbf{Departed Lux Industries
Inc.\ and took the CTO role at the Counterparty}, where he and a
substantially overlapping team rebuilt the Counterparty's ATS / BD /
TA stack from the ground up importing the LEL- and LRL-PR-tier Lux
modules (see \S IV(C), (F)). His operational authority covers the
platform rebuild and ongoing technical posture. \textbf{Mr.\ Kelling
is NOT named in any contemplated litigation by Lux Industries Inc.}
The officers with the affirmative duty to identify and remediate the
unlicensed-IP / OSS-licence-review failure operate \emph{above}
the Counterparty CTO role:
\begin{itemize}[leftmargin=2em,itemsep=0.1em]
\item Mr.\ Sonsurkar (CSO) bore the software-supply-chain-integrity
and third-party-IP-compliance duty by job function.
\item Mr.\ Trombley (CSO/Strategy, previously Acting CEO) bore
strategic and investor-representation responsibility for the
underlying technology stack.
\item Mr.\ Choi (CEO) bore ultimate disclosure responsibility for
representations to investors --- though, on information and belief,
he was \textbf{not informed} of the LEL / LRL-PR licensing
mismatch by Mr.\ Sonsurkar or Mr.\ Trombley despite their
respective duties, which mitigates his culpability from knowing
fraud toward an officer-level disclosure failure absent personal
knowledge.
\end{itemize}
The technical executor of the rebuild (Mr.\ Kelling) is not the
officer the litigation theory of the case names; the licensure-
review duty travels with the CSO seat and the disclosure duty
travels with the CEO/strategy chain. \textbf{Mr.\ Kelling's
continued tenure as Counterparty CTO is a Lux-required mandatory
term} in any partnership outcome (\S X, \S XI) --- the equity Lux
takes is rendered substantially worthless if the Counterparty
loses Mr.\ Kelling's technical leadership, because the current
Counterparty operational team without him is not viewed as
competent to maintain the platform. \textbf{See \S IV.E.}
\item \textbf{Eric Choi} --- Chief Executive Officer of the
Counterparty. Top operational authority; ultimate responsibility for
the Counterparty's representations to investors and regulators
regarding IP ownership and licensing.
\item \textbf{Austin Trombley} --- Chief Strategy Officer; previously
Acting Chief Executive Officer. Blockchain / AI engineer. Built a
\textbf{prior} (not current) ATS infrastructure at a prior employer
(Franklin Templeton) and \textbf{served as an eyewitness in the
Franklin Templeton litigation} arising from that engagement.
\textbf{Counsel should review the Franklin Templeton litigation
record} for (i) pattern-of-conduct evidence relevant to the present
matter, (ii) any limits on or adverse characterizations of his prior
testimony that could be impeachable, and (iii) the technical scope and
quality of the ATS he built there.

Material context: prior to Mr.\ Kelling's arrival as active Chief
Technology Officer at the Counterparty, the Counterparty's ATS as
built under Mr.\ Trombley's technical direction \textbf{lacked even
basic time-price priority order matching}, the foundational principle
of any functioning alternative trading system. The pre-Kelling
technical state is addressed in \S IV(F).
\item \textbf{Coleman Church} --- Chief Executive Officer of the
Counterparty's broker-dealer affiliate. FINRA-regulated responsibility
for the BD's technology and operational integrity.
\item \textbf{Jayant Sonsurkar} --- Chief Security Officer of the
Counterparty. Long-standing owner / admin on the Counterparty's GitHub
organizations: historically \texttt{satschel/*}, and currently
\texttt{liquidityio/*} --- the latter being the org under which the
rebuild integrating unlicensed Lux modules is being performed. The
CSO role at any regulated digital-asset venue carries the duty to
oversee software-supply-chain integrity, including open-source-license
review and third-party-IP compliance. As admin on the Counterparty's
own repositories, Mr.\ Sonsurkar had on-platform visibility into the
imported Lux LICENSE files --- LEL, LRL--PR, BSD-3, and others ---
as they entered the Counterparty's build tree. (For clarity:
\texttt{luxfi/*} GitHub-org admin was held by Mr.\ Nandy and Mr.\
Kelling on the Lux side; Mr.\ Sonsurkar did not have, and was not
required to have, admin access on Lux's own repositories. His
duty of OSS-license review attaches to imports into the
Counterparty's tree, where his admin access was present.)
\end{itemize}

\subsection*{Co-Plaintiffs / Co-Victims}

This memorandum is prepared on behalf of \textbf{Lux Industries Inc.}\
as primary plaintiff and \textbf{Hanzo, Inc.}\ as
\textbf{co-plaintiff / direct co-victim with independent and additive
claims}:

\begin{itemize}[leftmargin=1.5em,itemsep=0.1em]
\item \textbf{Hanzo, Inc.}\ is owned by substantially overlapping
investors / partners with Lux Industries Inc. Hanzo operates the
AI and infrastructure stack that is contributed back to the Lux
upstream as OSS, \textbf{and} is the \textbf{source-licensor of the
Hanzo-origin portion of the work product (the cryptographic
primitives library, the consensus-algorithm prototypes, the 2018
ESX-JV BD + ATS development, the Beyond Equity / Beyond Alpha
precedent exchange tech, and the early DEX work)} that was
\textbf{licensed / assigned forward into the Lux line of vehicles
for Lux's use} (see \S IV.A(1B)). \textbf{No such license or
assignment was ever made from Hanzo, Inc.\ to the Counterparty.}
Hanzo therefore holds its own independent and additive claims
against the Counterparty in respect of the Hanzo-source portion of
the contributed work product, separately pleadable as Counts I
(DTSA), II (Copyright), III (Patent), V (Aiding-and-Abetting), VI
(Unjust Enrichment), VII (Conversion), and XII (Civil Conspiracy)
in Hanzo, Inc.'s name. The Counterparty's exposure on the Hanzo-
source layer is independent of and additive to its exposure on the
Lux-source layer; any partnership-or-settlement construction
contemplates release of \textbf{both} Lux Industries Inc.\ and
Hanzo, Inc.\ on final execution (\S XI ¶21--22), or none.
Hanzo's forward operational requirements additionally include
\textbf{permanent access to the Counterparty's protocol and venue
for AI / infra-platform purposes} (\S XI ¶21).
\item \textbf{Zoo (Zoo Labs Foundation)} likewise requires
permanent access to the Counterparty's protocol and venue for AI-
research / decentralized-science purposes (see \S XI ¶21).
\item The Beyond Equity Markets and Beyond Alpha Ventures trading-
platform monorepos (referenced in \S IV.A(3) as the precedent
Alpaca-integrated exchange tech) are Lux / Hanzo work product,
not Counterparty work product; their inclusion in this memorandum
as evidence of pre-Counterparty exchange-tech precedent runs to
both Lux and Hanzo standing.
\item The \textbf{2018 ESX JV} BD + ATS development work product
(Hanzo, Inc.\ / Ryan Kavanaugh / Triller; Hanzo / Lux side received
a 20\% stake in ``Entertainment Stock X'') is documented Hanzo-origin
upstream that further supports Hanzo, Inc.'s direct co-plaintiff
standing under Counts I--III, V--VII, and XII.
\end{itemize}

For purposes of the substantive counts and damages analysis in
\S V--\S VI, ``Lux'' should be read to include Hanzo where the
context requires (i.e.\ for the OSS contribution chain-of-title,
for the Hanzo-source IP layer never licensed to the Counterparty,
and for the permanent-access mandatory term in \S XI ¶21).

\subsection*{Counterparty Entity Structure}

For purposes of this memorandum, ``the Counterparty'' refers
collectively to the operating entities under the current capital
raise (a Delaware corporation parent and its US LLC operating
affiliate, transacting as Liquidity.io) unless context requires
otherwise. The corporate structure should be confirmed during
outside-counsel diligence.

\textbf{Lead investor (current round).} Forum Markets, a publicly-
traded company (formerly Ethzilla; ticker \textbf{FRMM} (NASDAQ; CIK 0001690080; principal office Palm Beach, FL; SIC 6199)), continues in the lead-investor role
from the prior round and has signalled continued leadership of the
current \$75M raise at the \$1.1B valuation. Forum Markets' public-
issuer status means any decision to anchor a material investment in
the Counterparty is potentially material to Forum Markets' own
shareholders, with attendant Form 8-K considerations.

% ─────────────────────────────────────────────────────────────────
\section*{III.\ Jurisdiction, Venue, and Applicable Law}

\textbf{Federal causes of action} sound in: Defend Trade Secrets Act
(DTSA), 18 U.S.C.\ \S\S 1836--1839; Copyright Act, 17 U.S.C.\ \S\S
101 \emph{et seq.}; Patent Act, 35 U.S.C.\ \S\S 271, 281, 284, 285;
Lanham Act, 15 U.S.C.\ \S 1125(a); Securities Exchange Act of 1934,
15 U.S.C.\ \S 78j(b) (Rule 10b-5, 17 C.F.R.\ \S 240.10b-5).
\textbf{State and common-law claims} sound in Delaware (breach of
fiduciary duty, \emph{Aronson v.\ Lewis}, 473 A.2d 805 (Del.\ 1984);
aiding and abetting, \emph{Malpiede v.\ Townson}, 780 A.2d 1075
(Del.\ 2001)), and California (the likely Counterparty principal
place of business): common-law fraud, unjust enrichment, conversion,
tortious interference with prospective economic advantage, and civil
conspiracy.

\textbf{Forum recommendation:} Federal District Court for the
Northern District of California for the principal action, with
related Delaware Chancery proceeding for the corporate-governance
counts. Personal jurisdiction over the individual Counterparty
officers is established by their commercial activities directed at
Lux (a Delaware corporation) and their California residence.

\textbf{Litigation hold} --- Lux should issue a formal preservation
notice to the Counterparty and to each named individual within 30
days of any term-sheet delivery. Failure of preservation supports an
adverse-inference instruction at trial.

% ─────────────────────────────────────────────────────────────────
\section*{IV.\ Statement of Facts}

\textit{The following allegations are stated on information and
belief, based on the personal knowledge of the undersigned in his
dual operational role, supplemented by Lux's internal records.}

\subsection*{A. Background of Lux Industries Inc.\ and Its IP (2008--2026)}

\textbf{1.} \textbf{All IP at issue is owned by Lux Industries Inc.}
The IP was developed across a continuous corporate-identity lineage
beginning in 2008, with successive assignment of work product
forward through each corporate vehicle until consolidation at
present-day Lux Industries Inc.:

\begin{itemize}[leftmargin=2em,itemsep=0.1em]
\item \textbf{2008+} --- \textbf{Monoid LLC}: Mr.\ Kelling's solo
Advanced AI / ML / Blockchain research firm. Monoid carried the
foundational cryptographic-primitives research, the early machine-
learning systems work, and the earliest distributed-systems
engineering work that subsequently fed forward into the consensus and
matching-engine designs at issue. \textbf{[Outside counsel to enumerate
additional earlier research vehicles in the 2008--2016 window during
diligence.]}
\item \textbf{2016} --- \textbf{Hanzo, Inc.}\ (Delaware corporation)
incorporated, succeeding the prior \textbf{Crowdstart} venture
operated by Mr.\ Kelling. The Lux work continued under Hanzo, Inc.\ as
the research and engineering vehicle for the next phase --- the
cryptographic primitives library, the consensus-algorithm prototypes,
and early DEX work that became the Beyond Equity / Beyond Alpha
precedent exchange tech (\S IV.A(3)). \textbf{In 2018, Hanzo, Inc.\
pursued its own original BD + ATS licensure development through a
joint venture (the \emph{``ESX JV''}) with Ryan Kavanaugh and the
Triller group; the Hanzo / Lux side received a 20\% stake in the
JV vehicle (``Entertainment Stock X'')}, and the BD + ATS
development work product was assigned / licensed forward into the
Lux line of vehicles. The broader research stack that subsequently
became Lux-branded IP \textbf{was likewise assigned forward into the
Lux line}; Hanzo, Inc.\ retains a residual licensor / co-source
posture in respect of the contributed work, and Mr.\ Kelling is
currently compensated in material part through \textbf{Hanzo, Inc.\
consulting fees} (see \S IV.A(1B)--(1C) below).
\item \textbf{\emph{[period to be confirmed by counsel]}} --- IP
and team transitioned to \textbf{Lux Partners Limited}, an Isle of
Man (IOM) private company limited by shares, which carried the
research and early-stage product development through the next phase,
including the original DEX matching-engine and threshold-MPC work.
\item \textbf{Currently} --- IP and operating activity sit in
\textbf{Lux Industries Inc.}, the present-day Delaware corporation
that holds \textbf{all of the IP at issue} in this memorandum and
is the operative Counterparty-facing licensor of the LEL- and
LRL-PR-tier modules.
\end{itemize}
The corporate-identity chain is unbroken in respect of the
underlying IP (Monoid LLC + earlier research vehicles (2008+)
$\rightarrow$ Hanzo, Inc.\ (2016, prior Crowdstart venture)
$\rightarrow$ Lux Partners Limited (IOM) $\rightarrow$ Lux
Industries Inc.\ (currently)), and the operative Counterparty-facing
licensor in 2026 is Lux Industries Inc.\ alone. \textbf{[Outside-
counsel diligence to confirm: full enumeration of pre-Hanzo research
vehicles; the assignment instruments between Monoid LLC, the
additional pre-2016 vehicles, Hanzo, Inc., Lux Partners Limited
(IOM), and Lux Industries Inc.; existence of those instruments
establishes unbroken chain-of-title for the patent, copyright, and
trade-secret counts.]}

\textbf{1A.\ Mr.\ Kelling's pre-Monoid precedent context.} As
supporting context for the chain-of-attribution and originality
claims at \S IV.A(4) (patents) and the damages analysis at \S VI,
Mr.\ Kelling's professional history before Monoid LLC includes
relevant founding- / officer-level experience in early digital
securities and equity-financing technology: \textbf{Chief Technology
Officer of Crowdfunder.com}, an early equity-financing platform; and
prior work on \textbf{Ar.ca}, an early digital-securities platform /
DEX. This pre-Monoid record establishes the technical depth and
continuity that underwrites the originality and uninterrupted
attribution arguments in \S IV.A(3)--(4). \textit{Counsel may include
or omit at discretion based on diligence; it is preserved here for
completeness.}

\textbf{1B.\ The Hanzo-source-license layer, and why it does not
extend to the Counterparty.} A material portion of the work
catalogued in \S IV.A(3) was originally developed under Hanzo, Inc.\
(2016 onward) and was \textbf{licensed and / or assigned forward
into Lux} (Lux Partners Limited (IOM), then Lux Industries Inc.) for
Lux's operative use. \textbf{No such license or assignment was ever
made from Hanzo, Inc.\ to the Counterparty (Liquidity.io or any of
its affiliates).} The Counterparty therefore has \textbf{two
independent licensor lacunae} in respect of the modules at issue: it
lacks the operative LEL v1.2 / LRL--PR v1.0 license from Lux
Industries Inc.\ (the Lux-side gap that drives Counts I--III and
VI--VII of this memorandum), \textbf{and} it lacks any upstream
license from Hanzo, Inc.\ in respect of the Hanzo-origin portion of
the contributed work product. Hanzo, Inc.\ is therefore named as a
co-Plaintiff / co-victim at \S II (Parties and Dramatis Personae)
and the Counts that are pleadable on the Hanzo-source layer
(specifically Counts I (DTSA), II (Copyright), III (Patent), VI
(Unjust Enrichment), and the conspiracy-and-aiding-and-abetting
counts (V, XII)) run in Hanzo, Inc.'s name as well as Lux
Industries Inc.'s, in addition to the unitary Lux counts. The
Counterparty's exposure on the Hanzo-source layer is independent of
and additive to its exposure on the Lux-source layer; the
partnership construction in \S\S X--XI accordingly contemplates a
release of \textbf{both} Lux Industries Inc.\ and Hanzo, Inc.\ on
final execution (\S XI ¶21--22), or none.

\textbf{1C.\ Mr.\ Kelling's current Hanzo consulting fees as
evidentiary fact.} Mr.\ Kelling is presently compensated, in
material part, through \textbf{Hanzo, Inc.\ consulting fees}. This
arrangement is consistent with Mr.\ Kelling's continuing role as a
contributor and intellectual-property source whose work has flowed
forward into the Lux line of vehicles for Lux's use. \textbf{It is
not consistent with --- and provides no support for --- any claim that
Mr.\ Kelling carried a free-and-clear personal license from Hanzo,
Inc.\ or from any predecessor entity into the Counterparty.} The
modules and engineering output Mr.\ Kelling contributed into the
Counterparty's build during October--December 2025 and onward
remain subject to (i) the Lux licensor posture for the Lux-branded
modules, and (ii) the residual Hanzo, Inc.\ source-licensor posture
for the Hanzo-origin portion. This fact pattern materially
strengthens the Unjust Enrichment (Count VI), Conversion (Count
VII), and Aiding-and-Abetting (Count V) theories.

\textbf{1D.} The technology stack so developed includes consensus
algorithms, post-quantum cryptographic primitives, threshold and
MPC protocols, fully-homomorphic-encryption (FHE) precompile
architecture, virtual-machine substrates, EVM derivatives,
decentralized-exchange matching engines, and cross-chain messaging
protocols.

\textbf{2.} Lux's direct development investment in this stack exceeds
\$20 million in recorded cost from 2016 to date \textbf{[evidentiary
source: sworn declarations of corporate-finance officers
across Hanzo, Inc., Lux Partners Limited (IOM), and Lux
Industries Inc., consolidating recorded personnel + contractor
+ infrastructure cost across the lineage period --- to be
prepared by counsel in advance of filing]}. Fully-loaded personnel
cost over the same period is a multiple of that figure.

\textbf{3.} Chronological development milestones:

\begin{itemize}[leftmargin=2em,itemsep=0.1em]
\item \textbf{2016 -- 2017}: Founding under Hanzo, Inc. Foundational
consensus research; early Quasar prototypes; cryptographic
primitives library; \texttt{secp256k1}, \texttt{Ed25519},
\texttt{BLS12-381} reference implementations.

\textbf{Hanzo, Inc.\ client / advisory portfolio in this period
(IP-relevant to the present matter)}:

\begin{itemize}[leftmargin=2em,itemsep=0.05em]
\item \textbf{Unikrn / Unikoin Gold (2017)} --- early regulated
e-sports bookmaker; Hanzo client / technical advisor on the
\textbf{2017 Unikoin Gold ICO}, which raised approximately
\textbf{120,000 ETH}. Hanzo, Inc.\ contributed cryptographic,
smart-contract, and infrastructure work that subsequently fed
forward into the Lux line of vehicles.
\item \textbf{Ikigai Fund} --- Hanzo, Inc.\ client; technical
contribution that subsequently fed forward into Lux.
\item \textbf{Ar.ca} --- early digital-securities platform / DEX;
Hanzo, Inc.\ client and engineering counterparty; predecessor
technical work that informs the originality and chain-of-
attribution analysis for the present DEX / matching-engine
modules. Mr.\ Kelling additionally held CTO-level operating
involvement at Ar.ca (\S IV.A(1A)).
\end{itemize}

The aggregate IP output of these engagements --- the cryptographic
primitives, the early smart-contract and digital-securities
patterns, and the operational platform work --- \textbf{was
licensed / assigned forward into the Lux line of vehicles} on the
same chain-of-attribution basis as the ESX JV work product
(\S IV.A(1)). Hanzo, Inc.\ retains the residual source-licensor
posture in respect of those contributions, and \textbf{no parallel
licence was granted from Hanzo, Inc.\ to the Counterparty}
(\S IV.A(1B)).
\item \textbf{2018 -- 2020}: Active development of the planet-scale
DEX matching engine with 1ms finality. Initial EVM fork (LGPL).
Threshold-MPC research. \textbf{Hanzo, Inc.\ pursued original BD +
ATS licensure development through the \emph{ESX JV} with Ryan
Kavanaugh / Triller (Hanzo / Lux side received a 20\% stake in
``Entertainment Stock X''); the BD + ATS work product from that JV
was assigned forward into the Lux line of vehicles}. CDAX Limited
(IOM affiliate) operated under an \textbf{IOM Class 8
money-transmission licence} in the same period, providing the
operational compliance footprint for the IOM operations.
\item \textbf{2020 -- 2023}: Quasar consensus stabilization; Pulsar
and Magnetar Apache-2.0 threshold-PQ primitives; FHE prototypes.

\textbf{Luxembourg digital-securities + tokenized-commodity
platform (Creatrust fund-manager partnership)}: in the same
period, Lux (and Lux Partners Limited as predecessor) built and
operated a regulated digital-securities platform in
\textbf{Luxembourg} (CSSF-supervised footprint, \S VIII Historical
native support) in partnership with \textbf{Creatrust} as the
licensed Luxembourg fund-manager counterparty. The platform
covers (i) tokenized commodities --- \textbf{uranium, gold, and
silver}; (ii) tokenized fund-of-funds (FoF) structures; and (iii)
tokenized fund LP interests --- all issued on Lux's \textbf{ERC--
3643 (T--REX)}\footnote{ERC-3643 / T-REX is the Tokeny-stewarded
EVM standard for permissioned, regulator-friendly security
tokens; the standard is the de facto EU-compliant security-token
substrate.} compliant smart-contract suite. The \textbf{same
exact ERC-3643-compliant contract layer} Lux developed and deployed
for the Creatrust / Luxembourg digital-securities platform is the
underlying contract stack now embedded in the Counterparty's ATS
build (\S IV.C) --- which is independently dispositive of the
chain-of-attribution and unjust-enrichment counts because there is
\textbf{no plausible independent-derivation defence} for the
Counterparty's smart-contract layer when the same contracts (down
to function signatures and storage layout) are deployed and
diligenceable on Luxembourg-CSSF supervised infrastructure that
predates the Counterparty's build.
\item \textbf{2024}: Ringtail lattice-signature scheme; FIPS 203 /
204 / 205 verifier precompiles (ML-KEM, ML-DSA, SLH-DSA); P3Q
protocol; LEL v1.2 and LRL--PR v1.0 formalized; Warp messaging
protocol stabilization; FHE-based dark-pool matching.
\item \textbf{2025}: LX Suite DEX precompile family (Pool / Oracle /
Router / Hooks at LP-9010 onward) released following Uniswap v4's
hooks model; Corona Ring-LWE threshold scheme (successor to Ringtail);
LP-4200 precompile address standardization; multi-arch GPU acceleration
under the closed \texttt{\textasciitilde/work/luxcpp/*} tree.

\textbf{October--December 2025 (critical Counterparty-enabling work):}
Mr.\ Kelling personally introduced \textbf{Alpaca} (the FINRA-registered
broker-dealer that currently serves as the Counterparty's clearing BD)
to the Counterparty, and over October--December 2025 personally
developed (and integrated, with the Lux engineering team's substantial
contribution) the exchange-side technology that allows the Counterparty
to interoperate with Alpaca's clearing, custody, and securities-
movement APIs. \textbf{Precedent context}: the Alpaca-integrated
exchange-tech stack Mr.\ Kelling brought to the Counterparty is a
substantial reuse / port of the prior \textbf{Beyond Equity Markets}
(\texttt{\textasciitilde/work/beyondequity}) and \textbf{Beyond Alpha
Ventures} (\texttt{\textasciitilde/work/beyondalpha}) trading-
platform / unified-financial-platform monorepos built by the Lux /
Hanzo engineering team prior to the Counterparty engagement. The
Counterparty's current Alpaca-integration layer is therefore not
Counterparty-original work --- it is a port of pre-existing Lux /
Hanzo platform code, which materially extends the IP-attribution
analysis and the unjust-enrichment count. \textbf{This work is foundational to the Counterparty's
ATS having been able to launch at all in its current commercial form}:
absent Alpaca clearing and the Lux-built Alpaca-integration layer, the
Counterparty had no operational path to executing customer trades
under a SEC-registered clearing arrangement. The relevant Alpaca-side
diligence material (Live Broker-API partner forms, beneficial-
ownership certifications, DD checklists signed by Counterparty
officers) is preserved as evidentiary record at
\texttt{liquidity/exhibits/alpaca-liquidity-bd/} adjacent to this
memorandum. Damages attribution and partnership equity-share both
weight this Alpaca-integration contribution heavily; absence of any
contemporaneous compensating consideration to Lux for this work
further supports the unjust-enrichment and unjust-licensing
positions (Counts II, VI, IX).
\item \textbf{2026}: Brand-neutral cleanup; relicensing of Pulsar /
Magnetar under Apache-2.0; ongoing institutional partnership
development.
\end{itemize}

\textbf{4.} \textbf{Patent portfolio --- bifurcated between Lux
Industries Inc.\ and Hanzo, Inc.\ licensors.}

\textit{Lux Industries Inc.\ portfolio} --- holds or has filed for
patents on, without limitation: Quasar BFT consensus; sampler-bias
correction; Ringtail / Corona lattice threshold schemes; Unified VM
Execution Interface with Pluggable Backends; Post-Quantum State
Verification with Verkle Proofs; FHE precompile architecture;
GPU-accelerated DEX matching engine; AI-optimized bytecode
execution. Source whitepapers and patent-grade specifications
preserved at \texttt{\textasciitilde/work/lux/papers/} (LP-* series,
including \texttt{lp-020-quasar-consensus}, \texttt{lp-073-pulsar},
\texttt{lp-137-gpu-crypto-stack}, \texttt{lp-305-quantum-secure-
assets}, \texttt{lp-307-quantum-consensus-protocol},
\texttt{lp-308-quantum-amm}, and the \texttt{lux-ats-architecture}
+ \texttt{gpu-evm-whitepaper} families).

\textit{Hanzo, Inc.\ portfolio} --- holds or has filed for patents
on, without limitation: AI / agent runtime architecture (Hanzo
ACI, Agent SDK, AGI platform); model-context-protocol (MCP)
infrastructure; AI-chain (ACI) cryptographic primitives bridging
inference and consensus; Crowdstart-vintage commerce / equity-
financing patterns; defense-grade AI / cross-domain / assured-
networking primitives. Source whitepapers preserved at
\texttt{\textasciitilde/work/hanzo/papers/} (including the
\texttt{defense/} subtree, \texttt{hanzo-aci}, \texttt{hanzo-agent-
*}, \texttt{hanzo-agi-platform}, \texttt{hanzo-ai-chain},
\texttt{hanzo-consensus-ai}, \texttt{hanzo-cloud-infrastructure},
\texttt{hanzo-blockchain-platform}, and the \texttt{jin/papers}
subtree).

\textit{Private GPU / C++ code layer.} A material portion of the
GPU-acceleration code (\textbf{the \texttt{\textasciitilde/work/luxcpp/}
private tree} and the GPU-EVM kernels) is maintained as
\textbf{private, non-open-source} source code by Lux and is
patent-bearing. This layer is independently a trade-secret
(Count I --- DTSA) and is patent-protected (Count III --- Patent).
The Counterparty's commercial deployment of any binary or shared-
object built against this layer is independently actionable on
both counts.

\textit{Cross-portfolio standing.} Because patents in the \emph{Lux}
portfolio and patents in the \emph{Hanzo, Inc.}\ portfolio are both
implicated by the Counterparty's current stack (the Lux-branded
modules touching the Lux portfolio; the Hanzo-AI / ACI / agent /
defense substrate touching the Hanzo portfolio, including via the
Lux\,$\to$\,Hanzo upstream OSS contribution chain), \textbf{Count III
(Patent Infringement) is pleaded in the names of both Lux Industries
Inc.\ and Hanzo, Inc.\ as co-plaintiffs}.

\textit{Catalog summary} (full diligence-ready table at
\texttt{liquidity/review/patent\_catalog\_2026\_05\_25.md}, adjacent
to this memorandum):

\begin{itemize}[leftmargin=2em,itemsep=0.05em]
\item \textbf{Eighty-seven (87) patent-bearing inventions} cataloged
in total.
\item Lux Industries Inc.\ exclusive: \textbf{47}; Hanzo, Inc.\
exclusive: \textbf{32}; joint Lux + Hanzo: \textbf{8}.
\item Patent posture: \textbf{11 filed}; 51 filable; 21 defensive
publications; 7 trade-secret records.
\item Counterparty \textbf{probable read-ons}: \textbf{23
inventions} (across Quasar 3.0 consensus, GPU-accelerated EVM
(\texttt{cevm}), Warp 2.0 cross-chain finality, Block-STM parallel
execution, Lightspeed DEX matching engine, Quantum AMM +
perpetuals, Active Semantic Optimization, ML-DSA-65 / Groth16
validator binding, Lux Bridge protocol, and related primitives).
\item \textbf{Top ten inventions by damages exposure} (3.1$\times$
to 5.2$\times$ exposure multipliers on a comparable-transaction
basis) sit at the core of the Counterparty's stated ``1\,ms
finality'' and ``GPU-accelerated matching'' architecture; the
Quasar 3.0 consensus and the Lightspeed DEX matching engine are
the primary litigation anchors.
\item \textbf{Estimated aggregate Count III damages exposure}
(reasonable royalty floor under 35 U.S.C.\ \S 284, before willful-
infringement trebling and before any DTSA / Copyright add-ons):
\textbf{\$180M -- \$420M, base case \$280M}.
\end{itemize}

\textbf{[Outside counsel to verify USPTO filing status and PCT
extensions of each enumerated invention; perfect any provisional
filings to non-provisional within the standard 12-month window for
any invention in active commercial use; coordinate Copyright Office
registrations for Count II availability of statutory damages.]}

\subsection*{B. The Lux Licensing Tier Structure}

\textbf{5.} Lux publishes its software under a three-tier strategy:

\textbf{6.} \textbf{Tier 1 (BSD-3-Clause)}: commodity infrastructure,
free commercial use. \texttt{luxfi/node}, \texttt{luxfi/sdk},
\texttt{luxfi/cli}, \texttt{luxfi/api}, \texttt{luxfi/codec},
\texttt{luxfi/database}, \texttt{luxfi/p2p}, and approximately twenty-
five (25) other modules.

\textbf{7.} \textbf{Tier 2 (LEL v1.2)}: source-available, patent-
protected. Per LEL \S\S 1 and 3, commercial rights granted only to
``Authorized Networks'' --- the Lux Primary Network (NetworkID 1,
ChainID 96369), Lux testnets and devnets, and ``L1/L2/L3 chains
descending from the Lux Primary Network.''

\textbf{8.} Per LEL \S 4, forks and competing networks prohibited.
Per LEL \S 5, all patent rights reserved. Per LEL \S 8, license
terminates automatically on breach. Modules: \texttt{luxfi/consensus},
\texttt{luxfi/crypto}, \texttt{luxfi/mpc}, \texttt{luxfi/threshold},
\texttt{luxfi/ringtail}, \texttt{luxfi/corona}, \texttt{luxfi/fhe},
\texttt{luxfi/precompile}, \texttt{luxfi/captable},
\texttt{luxfi/broker}, \texttt{luxfi/treasury},
\texttt{luxfi/transfer}, \texttt{luxfi/compliance},
\texttt{luxfi/hsm}, \texttt{luxfi/aml}, \texttt{luxfi/maker},
\texttt{luxfi/warp}, among others.

\textbf{9.} \textbf{Tier 3 (LRL--PR v1.0)}: research-only license.
Section 3 expressly prohibits commercial use; Section 4 reserves all
patent claims. Modules: \texttt{luxfi/vm}, \texttt{luxfi/staking},
\texttt{luxfi/sampler}, \texttt{luxfi/dex} (Go bindings prior to its
2026 move to fully proprietary).

\textbf{10.} \textbf{Tier 4 (Closed Proprietary)}: the
\texttt{\textasciitilde/work/luxcpp/*} tree (CUDA / Metal / WebGPU /
FPGA / lattice GPU accelerators, the C++ DEX matching engine). Not
currently in the Counterparty's production graph; access requires a
separately negotiated Private-Moat Addendum.

\subsection*{C. The Counterparty's Commercial Use of Lux IP}

\textbf{11.} The Counterparty operates a registered Alternative
Trading System, broker-dealer, and transfer-agent business in the
United States, commercially deploying the validator binary
\texttt{lqd} (built from \texttt{liquidityio/node}) on three live
blockchain networks (Devnet ChainID 8675312, Testnet ChainID 8675310,
Mainnet ChainID 8675309), each of which incorporates dozens of Lux
modules.

\textbf{12.} The Counterparty further deploys the Liquidity EVM,
ATS, BD, and TA backend services, white-label tenants (MLC, VCC,
TZERO, ReTokens, and others), and approximately 12,796 SecurityTokens
on Devnet alone running on top of Lux precompiles.

\textbf{13.} The companion audit (referenced internally as the
license inventory) catalogues every Lux module directly imported.
\textbf{Twenty-three (23)} directly-imported modules carry LEL or LRL-PR tier licenses requiring an executed commercial license (verified count: \textbf{19 LEL + 4 LRL-PR}, enumerated against \texttt{go.mod} files across the Counterparty tree).

\textbf{14.} \textbf{No commercial license has ever been executed
between Lux Industries Inc.\ and the Counterparty.} Outside counsel
should verify the non-existence of any such instrument through
document discovery.

\textbf{15.} \textbf{The Liquid chains DO activate the Lux
Warp messaging precompile at the canonical address
\texttt{0x0200000000000000000000000000000000000005}}, as verifiable
in the Counterparty's own \texttt{configs/network/chains/C.json}
and \texttt{cchain-genesis.json} files: \texttt{warpConfig:
\{blockTimestamp:\,0,\,\dots\}} activates Warp at genesis, and
the precompile allocation at \texttt{0x02000\dots05} carries
\texttt{code:"0x01"}. \textbf{Warp activation strengthens, rather
than weakens, the LEL/LRL-PR licensing claim}: the Liquid chains
are not merely consuming Lux modules in source form; they are
operating as a derivative chain consuming the Lux interoperability
plane in production. \textbf{Under the LEL Authorized-Network
definition (LEL v1.2 \S 1, \S 3, \S 5), the operative test is
the \emph{Descending-Chain} criterion --- whether the chain
descends from Lux consensus, retains the Lux-defined precompile
namespace, and remains within the Lux interoperability domain. The
Liquid chains satisfy the Descending-Chain criterion in code and
therefore fall within the LEL grant's commercial-use trigger,
which requires an executed commercial license that has never been
executed (see \textbf{14.}).}

\subsection*{D. The Counterparty Officer Team and Knowledge}

\textbf{16.} Mr.\ Eric Choi (CEO) holds ultimate operational
responsibility for technology licensing and IP compliance. His
signature appears (or will appear) on the subscription documents for
the current \$75M / \$1.1B-valuation raise; representations therein
sound in his name. As of the date of this memorandum the round has
not closed; the duty of disclosure is therefore current and ongoing.
\textbf{Mitigating context}: on information and belief, Mr.\ Choi
was \textbf{not informed} of the LEL / LRL-PR licensing mismatch by
Mr.\ Sonsurkar (CSO) or Mr.\ Trombley (CSO/Strategy) despite the
affirmative duties of those officers to flag it. Mr.\ Choi's
culpability is accordingly mitigated from knowing fraud toward an
officer-level disclosure failure absent personal knowledge, which
materially shifts the Count IX (10b-5) and Count XI (common-law
fraud) analysis as to him.

\textbf{17.} Mr.\ Austin Trombley (CSO, prev.\ Acting CEO) ---
discussed in \S II above, including the Franklin Templeton litigation
history and the pre-Kelling ATS technical state.

\textbf{18.} Mr.\ Coleman Church (BD CEO) --- FINRA-regulated
responsibility for the BD's technology integrity; representations on
FINRA filings sound in his name.

\textbf{19.} Mr.\ Jayant Sonsurkar (CSO) --- Long-standing owner /
admin on the Counterparty's GitHub organizations \texttt{satschel/*}
(historic) and \texttt{liquidityio/*} (current rebuild org), giving
him direct, contemporaneous, ongoing visibility into every Lux
\texttt{LICENSE} file \textbf{as those files entered the
Counterparty's build tree} on import --- the LEL, LRL-PR, BSD-3,
GPL, LGPL, BUSL, MIT, and Apache markers travel with their source
files into any tree that consumes them. \texttt{luxfi/*} GitHub-org
admin was held by Lux's own engineering team (Mr.\ Nandy, Mr.\ Kelling)
on the upstream side; Mr.\ Sonsurkar did not have admin on
\texttt{luxfi/*} itself, nor was he required to --- his
software-supply-chain integrity duty attaches to imports into the
Counterparty's own tree, where his admin access was present.

\textbf{20.} The CSO role at a regulated digital-asset venue
carrying ATS / BD / TA registrations includes (without limitation)
software-supply-chain integrity oversight, identification of LEL /
LRL--PR / GPL / LGPL restricted licenses in deployed dependencies,
escalation of licensing mismatches to the CEO / BD CEO / GC, and
implementation of remediation.

\textbf{21.} On information and belief, Mr.\ Sonsurkar did not, at
any time prior to the elevation of this matter by Mr.\ Nandy, raise the
LEL / LRL--PR mismatch internally, request a commercial license,
include the unlicensed-IP status in regulatory compliance reviews, or
flag the issue during current-raise diligence. By contrast, Mr.\ Nandy
(VPE at the Counterparty, concurrently CTO at Lux) did identify the
mismatch through his own dual-role visibility and elevated it
upstream. The omission by Mr.\ Sonsurkar --- the officer to whom that
identification properly belonged --- is the central evidentiary fact
supporting Count IV (breach of fiduciary duty) and Count V (aiding
and abetting).

\textbf{22.} Messrs.\ Choi, Trombley, and Church each had operational
or strategic visibility into the Counterparty's technology stack
during the relevant period; the precise scope of each officer's
knowledge will require diligence development.

\subsection*{E. The Position of Mr.\ Zach Kelling}

\textbf{23.} Mr.\ Kelling is the Founder of both Lux Industries
Inc.\ and the Counterparty. He \textbf{has departed Lux Industries
Inc.}\ and is \textbf{currently Chief Technology Officer of the
Counterparty (Liquidity.io)}. His operational authority at the
Counterparty covers the platform rebuild and ongoing technical
posture. The Counterparty's commercial decisions and investor-facing
representations remain with Mr.\ Choi (CEO) and Mr.\ Trombley (CSO).

\textbf{24.} As a former Lux Founder/officer/cryptographer who
departed Lux Industries Inc.\ and took the Counterparty CTO role
during which the Lux-origin IP rebuild occurred, Mr.\ Kelling is a
\textbf{prospective named defendant} in the available causes of
action --- specifically Count I (Misappropriation of Trade Secrets,
DTSA), Count II (Copyright Infringement), Count III (Patent
Infringement), Count IV (Breach of Fiduciary Duty owed to Lux as
former Lux fiduciary), Count VI (Unjust Enrichment), Count VII
(Conversion), Count IX (Rule 10b-5 Securities Fraud, by virtue of his
current CTO role during the active \$75M / \$1.1B-valuation raise),
Count XI (Common-Law Fraud), and Count XII (Civil Conspiracy).
Mr.\ Kelling's Counterparty equity position is to be confirmed in
diligence; his former-Lux minority-shareholder position is
documented.

\textbf{25.} \textbf{Key-person retention --- mandatory in any
partnership outcome.} Despite Mr.\ Kelling's prospective defendant
posture in litigation, the partnership outcome (\S X) requires
him to \textbf{remain in place as Counterparty CTO post-settlement}.
Lux's equity stake in the Counterparty is rendered substantially
worthless without Mr.\ Kelling's continued technical leadership ---
the current Counterparty operational team, absent Mr.\ Kelling, is
not viewed as competent to maintain or extend the rebuilt platform
that underpins the Counterparty's enterprise value. The partnership
deal therefore pairs:
\begin{itemize}[leftmargin=2em,itemsep=0.1em]
\item \textbf{Mutual general releases of Mr.\ Kelling individually}
from all Counts in this memorandum (conditional on full execution
of the partnership agreement and ongoing performance);
\item \textbf{A multi-year key-person retention commitment} from the
Counterparty to retain Mr.\ Kelling as CTO with defined good-leaver
/ bad-leaver triggers, garden-leave constraints, and an equity
clawback or compensation structure tied to the partnership equity
value;
\item \textbf{Carve-outs} of Mr.\ Kelling from any litigation Lux
brings against other Counterparty officers (Choi / Trombley /
Church / Sonsurkar) if partnership outcome is achieved.
\end{itemize}
Without partnership, Mr.\ Kelling stands as a co-defendant on the
Counts above and the releases are not available. \textbf{See \S X for
the partnership tier structure and \S XI for the full mandatory-terms
list including key-person retention.}

\subsection*{F. Pre-Kelling Technical State and the ACH Exploit}

\textbf{26.} In the period immediately preceding Mr.\ Kelling's
arrival as active Chief Technology Officer, the Counterparty's ATS
platform was materially non-functional as a regulated alternative
trading system. Specifically, on information and belief:

\begin{itemize}[leftmargin=2em,itemsep=0.1em]
\item The ATS \textbf{did not implement working time-price priority
order matching} --- the foundational ordering principle of every
regulated US alternative trading system, required by SEC Reg ATS
for any venue purporting to operate as a securities trading facility.
\item No functioning ATS pipeline existed end-to-end.
\item \textbf{One day prior to Mr.\ Kelling's arrival as active CTO,
the platform suffered an ACH-funding exploit that resulted in the
theft of approximately \$40,000} via a trivial attack vector that
should have been remediated by even baseline security review. The
exploit evidences both the security posture of the platform under
the pre-Kelling officer team and the absence of competent CSO
oversight.
\end{itemize}

\textbf{27.} Upon his arrival, Mr.\ Kelling and his team (substantially
overlapping with the Lux engineering team) rebuilt the ATS / BD / TA
stack from the ground up, importing the Lux modules described in
\S IV(C) as the foundation. \textbf{In substantial part, the
technology that supports the current \$75M / \$1.1B-valuation raise was the Lux
IP that Mr.\ Kelling brought across, integrated with the regulatory
wrapper.}

\textbf{28.} The factual proposition that the Counterparty's
enterprise value as represented to investors is materially attributable
to Lux-origin technology contributed by Mr.\ Kelling and the broader
Lux team is supported by the contemporaneous technical record. This
context bears on damages attribution (\S VI), Mr.\ Trombley's
credibility as a technology witness, the breach-of-fiduciary-duty
count, and the protective treatment of Mr.\ Kelling (\S XII).

\subsection*{G. The Current \$75M / \$1.1B-Valuation Capital Raise (Forum Markets, Lead)}

\textbf{29.} The Counterparty is \textbf{currently undertaking a
capital raise of approximately \$75 million at a \$1.1 billion
valuation}. The round is led by \textbf{Forum Markets} (a
publicly-traded company, formerly Ethzilla; ticker \textbf{FRMM} (NASDAQ; CIK 0001690080; principal office Palm Beach, FL; SIC 6199)), continuing in the lead-investor
role from the prior round. As of the date of this memorandum, the
Counterparty has \textbf{verbal commitments covering at least
approximately half of the \$75 million round} (i.e.\ approximately
\$37.5M soft-circled). \textbf{The round has not closed.} Funds have
not been wired against current-round subscriptions; subscribers have
not been issued securities in respect of the current round; the
Counterparty's current and ongoing duty of disclosure to those
subscribers remains live until closing.
The subscription agreements, PPM, and side letters require outside-
counsel review for representations regarding (a) Counterparty
ownership of the underlying technology stack, (b) licensing status of
all third-party IP, (c) related-party transactions, (d) the CSO's
and BD CEO's compliance attestations, and (e) technology disclosures
in the risk factors section.

\textbf{30.} To the extent the current-raise subscription documents represent or omit
material facts regarding the Counterparty's ownership or licensed
access to Lux IP, those representations are independently actionable
under Section 10(b) of the Securities Exchange Act and Rule 10b-5.
See Count IX, infra.

% ─────────────────────────────────────────────────────────────────
\section*{V.\ Causes of Action Available to Lux Industries Inc.}

\subsection*{Count I --- Misappropriation of Trade Secrets (DTSA, 18 U.S.C.\ \S 1836) --- Bifurcated Standing (Lux + Hanzo)}

\textit{Lux Industries Inc.\ as plaintiff} --- Lux's unpublished
optimizations, integration know-how, architectural documents,
deployment recipes, MPC operational procedures, sampler-bias
remediation, and \textbf{pre-publication private-repository code
including the \texttt{\textasciitilde/work/luxcpp/} GPU layer}
qualify as trade secrets.

\textit{Hanzo, Inc.\ as co-plaintiff} --- Hanzo's unpublished AI /
agent runtime know-how, the ACI cryptographic primitives, the
defense-grade primitives, the Crowdstart-vintage commerce / equity-
financing patterns, and the residual unpublished portions of the
upstream OSS contribution chain that flowed into the Lux modules
likewise qualify as Hanzo trade secrets; Hanzo's standing on
Count I is independent of and additive to Lux's.

\textbf{Both Mr.\ Trombley AND Mr.\ Sonsurkar} hold continuing
owner / admin access to the Counterparty's GitHub organizations
(\texttt{satschel/*} and \texttt{liquidityio/*}), giving them
\textbf{100\% oversight and overview of every repository, every
\texttt{LICENSE} file as those files entered the Counterparty's
build tree, every commit, and every pull-request discussion}.
Their continuing repository access without an executed Lux \emph{or
Hanzo} commercial licence, coupled with the licensing-mismatch
silence on both sides of their respective duties (CSO software-
supply-chain integrity / OSS-licence review for Mr.\ Sonsurkar;
CSO/Strategy + prior Acting-CEO investor-representation
responsibility for Mr.\ Trombley), supports misappropriation by
improper means.

\textbf{Remedies (joint)}: Injunction (18 U.S.C.\ \S 1836(b)(3)(A));
actual damages and unjust enrichment (\S (B)); reasonable royalty
in lieu (id.); exemplary damages up to 2$\times$ for willful and
malicious misappropriation (\S (C)); attorneys' fees (\S (D));
seizure if record warrants (\S 1836(b)(2)). Damages run separately
to each co-plaintiff in respect of the trade secrets in its own
record that are misappropriated.

\subsection*{Count II --- Copyright Infringement (17 U.S.C.\ \S 501) --- Bifurcated Standing (Lux + Hanzo)}

\textit{Lux Industries Inc.\ as plaintiff} --- Lux owns copyright
in every source file under \texttt{luxfi/*} (public open-source
modules + the private \texttt{\textasciitilde/work/luxcpp/}
GPU layer). Copyright notices are present in source headers.

\textit{Hanzo, Inc.\ as co-plaintiff} --- Hanzo, Inc.\ owns
copyright in every source file under \texttt{hanzoai/*} (and
predecessor / sibling orgs reflecting earlier Hanzo client work
including \texttt{ar.ca}, the \texttt{unikrn} / \texttt{unikoin-
gold} engagement record, and the \texttt{crowdstart} commerce
work product) to the extent any such source file or derivative
work was reused or ported into the Counterparty's build tree.

The Counterparty has copied by direct import in \texttt{go.mod},
container-image embedding, binary distribution to validator nodes,
and on-chain deployment of derivative SecurityTokens. Copying that
reads on a \texttt{luxfi/*} header runs to Lux Industries Inc.;
copying that reads on a \texttt{hanzoai/*} (or predecessor)
header runs to Hanzo, Inc.

\textbf{Remedies (joint)}: Actual damages and infringer's profits,
\S 504(b); or statutory damages up to \$150,000 per work willfully
infringed, \S 504(c)(2); injunction (\S 502); impoundment (\S 503);
attorneys' fees and costs (\S 505). Statutory-damages aggregation
runs separately to each co-plaintiff in respect of its own
registered or registrable works. \textbf{[Outside counsel: confirm
Copyright Office registration status of each module under both
luxfi and hanzoai header; perfect any unregistered works prior to
filing to preserve statutory-damages availability.]}

\subsection*{Count III --- Patent Infringement (35 U.S.C.\ \S 271) --- Bifurcated Standing (Lux Industries Inc.\ + Hanzo, Inc.)}

\textit{Lux Industries Inc.\ as plaintiff} --- Lux holds or has
filed for patents covering the inventions enumerated in
\S IV(A)(4) (Lux portfolio). The Counterparty's deployment of
Quasar consensus, Ringtail / Corona thresholds, the sampler-bias
correction, the Liquid EVM substrate, the FHE precompile
activations, the DEX matching engine, and any binary built against
the private \texttt{\textasciitilde/work/luxcpp/} GPU layer each
implicate one or more patent claims in the Lux portfolio.

\textit{Hanzo, Inc.\ as co-plaintiff} --- Hanzo, Inc.\ holds or
has filed for patents covering the AI / agent runtime, MCP, ACI,
Crowdstart-vintage commerce, and defense-grade AI / cross-domain /
assured-networking primitives enumerated in \S IV(A)(4) (Hanzo
portfolio). To the extent the Counterparty's stack embeds or reads
on these primitives (whether directly or via Lux upstream OSS
contribution), Hanzo, Inc.\ holds independent and additive standing
to assert Count III against the Counterparty.

\textbf{Remedies (joint)}: Reasonable royalty floor (35 U.S.C.\
\S 284); lost profits where established; injunction (\S 283); treble
damages for willfulness (\S 284); attorneys' fees in exceptional
cases (\S 285). Damages run separately to each co-plaintiff in
respect of the patents in its own portfolio that are infringed.

\subsection*{Count IV --- Breach of Fiduciary Duty (Del., against Jayant Sonsurkar individually)}

As a corporate officer, defendant owes the corporation duties of
care and loyalty. Mr.\ Sonsurkar's failure to identify and flag the
LEL / LRL--PR licensing mismatch throughout his tenure is at minimum
gross negligence in the CSO duty of care; if the record supports
knowledge or willful blindness, it crosses into duty-of-loyalty
breach.

\subsection*{Count V --- Aiding and Abetting Breach (Del., against Messrs.\ Choi, Trombley, Church)}

Per \emph{Malpiede v.\ Townson}, 780 A.2d 1075 (Del.\ 2001): (i) a
fiduciary relationship; (ii) breach; (iii) knowing participation by
the non-fiduciary; (iv) damages. To the extent the named officers
had actual or constructive knowledge of the mismatch and continued to
direct commercial operations, they aided and abetted Mr.\ Sonsurkar's
breach.

\subsection*{Count VI --- Unjust Enrichment --- Bifurcated Standing (Lux + Hanzo)}

The Counterparty's commercial deployment of Lux \emph{and Hanzo}
IP without an executed license enriches the Counterparty at the
expense of each upstream licensor; the LEL and LRL--PR (Lux) and
the residual Hanzo source-licence posture (\S IV.A(1B)) each
expressly negate any justification. Standing runs separately to
each co-plaintiff in respect of its own contributed work product.

\subsection*{Count VII --- Conversion --- Bifurcated Standing (Lux + Hanzo)}

The Counterparty's deployment of LRL--PR-protected patent-bearing
Lux modules constitutes intentional interference with Lux's
exclusive rights; deployment of the Hanzo-source AI / agent / ACI
primitives constitutes parallel interference with Hanzo's
exclusive rights. Each interference is independently actionable
and runs to its respective licensor.

\subsection*{Count VIII --- Tortious Interference with Prospective Economic Advantage}

The Counterparty's white-label tenants (MLC, VCC, TZERO, ReTokens)
constitute particular instances of interference with Lux's
commercial-license pipeline with third-party ecosystem partners.

\subsection*{Count IX --- Securities Fraud (Rule 10b-5)}

To the extent the current \$75M / \$1.1B-valuation raise is supported by representations or
omissions regarding the Counterparty's IP ownership / licensing
status, elements of Rule 10b-5 are met. \textbf{Lux's primary
utility on this count is collateral leverage rather than direct
recovery; standing analysis is for outside counsel.}

\subsection*{Count X --- Lanham Act False Advertising (15 U.S.C.\ \S 1125(a))}

The Counterparty's marketing materials describe its technology stack
(post-quantum cryptography, FHE-based dark-pool matching, MPC
custody, multi-chain architecture, Quasar-derived consensus) in
terms that ascribe ownership or proprietary control without
disclosing the Lux origin or the conditional / unlicensed nature of
the access.

\subsection*{Count XI --- Common-Law Fraud}

To the extent Counterparty officers made representations to Lux
personnel that licensing arrangements would be formalized in due
course, and Lux relied to its detriment by continuing engineering
contribution without an executed agreement, common-law fraud lies.

\subsection*{Count XII --- Civil Conspiracy}

The Counterparty officer team's concerted commercial deployment of
unlicensed Lux IP and capital-raising on the back of that deployment
supports civil-conspiracy joint and several liability.

% ─────────────────────────────────────────────────────────────────
\section*{VI.\ Damages Analysis}

\subsection*{Small Scenario (Path A --- Negligent Posture)}

\begin{center}
\begin{tabular}{lr}
\toprule
Back-license value (12 yr at \$2.5M/yr base) & \$30M \\
Pre-judgment interest (5\%, simple) & \$10M \\
Royalty component (1.5\% of estimated GR) & \$10M \\
Forward-looking license value (NPV) & \$30M \\
\midrule
\textbf{Subtotal --- Small Scenario} & \textbf{\$80M} \\
\bottomrule
\end{tabular}
\end{center}

\subsection*{Medium Scenario (Path B --- Mixed-Knowledge, Contested)}

\begin{center}
\begin{tabular}{lr}
\toprule
Small-Scenario base & \$80M \\
DTSA 2$\times$ exemplary on willful misappropriation & +\$30M \\
Patent willfulness component (selective trebling) & +\$40M \\
Disgorgement on current-raise attribution (5\% of enterprise value, TODO: rebase per counsel) & +\$55M \\
Attorneys' fees and costs & +\$10M \\
\midrule
\textbf{Subtotal --- Medium Scenario} & \textbf{\$215M} \\
\bottomrule
\end{tabular}
\end{center}

\subsection*{Big Scenario (Path B Fully Litigated, Adverse Officer Record)}

\begin{center}
\begin{tabular}{lr}
\toprule
Base + interest + royalty & \$80M \\
DTSA full 2$\times$ exemplary & +\$80M \\
Patent full 3$\times$ trebling on willfulness & +\$120M \\
Disgorgement on current-raise attribution (15--20\% of enterprise value, TODO: rebase per counsel) & +\$200M \\
Lanham Act + Lanham fees & +\$10M \\
Securities fraud (if Lux can establish standing) & +\$50M \\
Attorneys' fees, costs, pre-judgment interest & +\$50M \\
\midrule
\textbf{Subtotal --- Big Scenario} & \textbf{\$590M} \\
\bottomrule
\end{tabular}
\end{center}

\subsection*{Opportunity-Cost Damages: 90-Day Bandwidth Diversion}

\textbf{The damages framework above does not yet capture a distinct
and significant category of harm: the opportunity cost to Lux of
having Mr.\ Kelling's --- and the broader Lux engineering team's ---
bandwidth diverted to the Counterparty's platform and to a
constellation of related ventures during a critical 90-day window.
The git record is unambiguous.}

\textbf{Quantification (rolling 90-day window) --- Counterparty
primary codebase.} Audit of the \texttt{liquidityio/*} GitHub
organization for commits attributable to Mr.\ Kelling:

\begin{center}
\begin{small}
\begin{tabular}{lrrr}
\toprule
\textbf{Repo} & \textbf{Zach commits} & \textbf{Total commits} & \textbf{Zach \%} \\
\midrule
\texttt{liquidityio/ats}        & 806    & 1{,}426 & 57\% \\
\texttt{liquidityio/bd}         & 388    & 751     & 52\% \\
\texttt{liquidityio/ta}         & 208    & 364     & 57\% \\
\texttt{liquidityio/mpc}        & 31     & 95      & 33\% \\
\texttt{liquidityio/evm}        & 23     & 32      & 72\% \\
\texttt{liquidityio/fhe}        & 10     & 14      & 71\% \\
\texttt{liquidityio/iam}        & 27     & 29      & 93\% \\
\texttt{liquidityio/node}       & 338    & 380     & 89\% \\
\texttt{liquidityio/operator}   & 354    & 406     & 87\% \\
\texttt{liquidityio/universe}   & 1{,}436 & 1{,}920 & 75\% \\
\texttt{liquidityio/contracts}  & 72     & 176     & 41\% \\
\texttt{liquidityio/dex}        & 24     & 32      & 75\% \\
\texttt{liquidityio/platform}   & 36     & 51      & 71\% \\
\texttt{liquidityio/superadmin} & 243    & 338     & 72\% \\
\texttt{liquidityio/exchange}   & 522    & 1{,}170 & 45\% \\
\texttt{liquidityio/console}    & 28     & 38      & 74\% \\
\texttt{liquidityio/ui}         & 27     & 143     & 19\% \\
\texttt{liquidityio/sdk}        & 30     & 30      & 100\% \\
\texttt{liquidityio/sdk-go}     & 9      & 9       & 100\% \\
\texttt{liquidityio/cli}        & 114    & 117     & 97\% \\
\midrule
\textbf{Subtotal (Counterparty primary, 90-day)} & \textbf{4{,}726} & \textbf{7{,}521} & \textbf{63\%} \\
\bottomrule
\end{tabular}
\end{small}
\end{center}

\textbf{Quantification --- related-venture codebases.} The audit
extends across five additional venture trees on which Mr.\ Kelling's
bandwidth was concurrently directed during the same 90-day window:

\begin{center}
\begin{small}
\begin{tabular}{lrrl}
\toprule
\textbf{Venture tree} & \textbf{Zach commits} & \textbf{LOC (add+del)} & \textbf{Note} \\
\midrule
VCC (vcc/exchange + contracts + wallet)        & 143 & 3{,}038{,}420 & \textbf{100\% new venture} (white-label tenant) \\
MLC (mlc/exchange + contracts + wallet)        & 153 & 3{,}044{,}377 & \textbf{100\% new venture} (white-label tenant) \\
EquityTable (loan-funding-engine + app)        & 98  & 681{,}666     & \textbf{100\% new venture} \\
OnyxPlus (admin + verify + onyxd + onboarding) & 188 & 57{,}379      & \textbf{Largely 100\% new venture} (KYC/ID JV) \\
LiquidityArchive (legacy + simplici-*)         & 683 & 9{,}541{,}566 & Legacy / archive maintenance \\
\midrule
\textbf{Subtotal (related ventures, 90-day)}   & \textbf{1{,}265} & \textbf{16{,}363{,}408} & \\
\bottomrule
\end{tabular}
\end{small}
\end{center}

\textbf{Grand totals across all six venture trees (90-day window):}

\begin{center}
\begin{tabular}{lr}
\toprule
\textbf{Total Zach commits (all 6 trees)} & \textbf{5{,}991} \\
\textbf{Total LOC churn (add + del across all 6 trees)} & \textbf{$\sim$93 million} \\
\textbf{LOC churn in Counterparty primary (top 8 repos)} & \textbf{$\sim$77 million} \\
\textbf{Distinct working days in \texttt{ats/} alone} & \textbf{55 (six days/week)} \\
\textbf{Repos where Zach $>$70\% of all commits} & \textbf{12 of 20 Counterparty repos} \\
\textbf{Repos where Zach $=$ 100\% of all commits} & \textbf{2 (\texttt{sdk}, \texttt{sdk-go})} \\
\bottomrule
\end{tabular}
\end{center}

\textbf{Three of the five related ventures (VCC, MLC, EquityTable) were
authored entirely from scratch during the 90-day window} ---
representing approximately \textbf{6.8 million lines} of net new code
across companies that the Counterparty offers to its tenants and
partners. A fourth (OnyxPlus, the KYC/ID joint venture) was largely
new in the same window. This is not maintenance contribution; this
is greenfield construction of multiple companies.

\textbf{Mr.\ Kelling's contribution to Lux upstream in the same
90-day window (the foundational Lux IP that powers the
Counterparty's stack):}

\begin{center}
\begin{small}
\begin{tabular}{lrr}
\toprule
\textbf{Lux repo} & \textbf{Zach commits} & \textbf{Total commits} \\
\midrule
\texttt{luxfi/node}             & \textbf{0} & 179 \\
\texttt{luxfi/consensus}        & \textbf{0} & 104 \\
\texttt{luxfi/crypto}           & \textbf{0} & 64 \\
\texttt{luxfi/evm}              & \textbf{0} & 99 \\
\texttt{luxfi/vm}               & \textbf{0} & 27 \\
\texttt{luxfi/precompile}       & \textbf{0} & 98 \\
\midrule
\textbf{Total (90-day window)} & \textbf{0} & \textbf{571} \\
\bottomrule
\end{tabular}
\end{small}
\end{center}

\textbf{The contrast is decisive}: 5,991 commits and $\sim$93M LOC
of Zach's bandwidth into Counterparty-side ventures versus exactly
zero commits to the upstream Lux repositories he is Founder and
Chief Cryptographer of.

\textbf{Value-creation rate translation.} The Counterparty's reported
current-round valuation in the active raise is approximately \$1.1 billion. The platform supporting that raise was substantially
built or rebuilt by Mr.\ Kelling during the 90-day window. The
implied value-creation rate is approximately \textbf{\$1.1B / 90
days = \$12.2M per day} of diverted attention --- and that figure
captures only the Counterparty's primary codebase; it does not yet
include the enterprise value of the three additional 100\%-new
ventures (VCC, MLC, EquityTable) and the largely-new OnyxPlus JV.
Counting those, the per-day value-creation rate of Mr.\ Kelling's
bandwidth during the 90-day window is materially higher than the
\$12.2M figure above.

\textbf{The opportunity cost to Lux} --- measured by what Lux's own
roadmap would have produced during the same window had Mr.\
Kelling's attention been available to Lux upstream --- \textbf{is
bounded above by that rate and below by the proportional founder-
share of Lux's enterprise trajectory}.

\begin{center}
\begin{tabular}{lr}
\toprule
\textbf{Component} & \textbf{Range} \\
\midrule
Bandwidth diversion at 50\% rate (90 days) & \$500M -- \$600M \\
At 25\% rate & \$250M -- \$300M \\
At 10\% rate (conservative) & \$100M -- \$120M \\
\midrule
\textbf{Conservative opportunity-cost claim} & \textbf{\$100--300M} \\
\bottomrule
\end{tabular}
\end{center}

This component is independent of and additive to the IP damages
above. \textbf{It is the single most striking quantitative fact in
the memorandum} and should be lead-in evidence for Mr.\ Dupont in
opening negotiation.

\subsection*{Damages Summary (Including Opportunity Cost)}

\begin{center}
\begin{tabular}{lrr}
\toprule
\textbf{Scenario} & \textbf{IP damages} & \textbf{+ Opp.\ Cost} \\
\midrule
Small (Path A, settlement context) & \$80M & \$180--380M \\
Medium (Path B, mixed-record litigation) & \$215M & \$315--515M \\
Big (Path B, full-record litigation w/ multipliers) & \$590M & \$690--890M \\
\bottomrule
\end{tabular}
\end{center}

\textbf{Partnership remedy (headline).} The figures above are
\textbf{litigation-cash damage ranges} --- the floor of what Lux could
recover in court. Lux's actual ask is a \textbf{partnership
agreement} structured as a \textbf{10\%--50\% equity issuance in the
Counterparty} plus a cash component plus reciprocal commercial
commitments. At the \textbf{partnership floor (10\% equity), nominal
value is \$120--135M} (\$110M equity at \$1.1B-valuation pricing +
\$10--25M cash). Recommended (25\%) is \$300--325M. Stretch (35\%) is
\$435--460M. Maximum (50\%) is \$600--650M. \textbf{Each tier carries
escalating governance rights} (board observer $\to$ seat $\to$ parity)
and broader reciprocal commitments (LEL grant $\to$ BaaS access $\to$
joint GTM $\to$ international distribution + closed-source GPU
access). The partnership remedy dominates the litigation-cash damages
because (i) it regularises the IP relationship directly, (ii)
preserves Counterparty operational continuity which is a stated
objective of Mr.\ Kelling (Counterparty CTO) and Mr.\ Nandy (Lux CTO /
Counterparty VPE), (iii) prices off the same \$1.1B valuation Forum
Markets and current-round investors are pricing into, and (iv) creates
ongoing economic alignment between Lux and the Counterparty for the
forward enterprise value Mr.\ Nandy's team will continue building.
See \S X for the Partnership Triangulation matrix.

% ─────────────────────────────────────────────────────────────────
\section*{VII.\ Trading-Fee Revenue Projections (User-Scale Analysis)}

The following projects the Counterparty's enterprise-value trajectory
at four user-base milestones and computes Lux's economic interest at
each under a proposed 10\% equity stake.

\subsection*{Annual Revenue, Enterprise Value, and Lux 10\% Stake}

\begin{center}
\begin{small}
\begin{tabular}{lrrrr}
\toprule
\textbf{Scenario} & \textbf{Users} & \textbf{Revenue} & \textbf{Enterprise Value} & \textbf{Lux 10\%} \\
\midrule
\multicolumn{5}{l}{\textit{Conservative ARPU (\$50/yr)}} \\
Stage 1 & 100K & \$5M & \$50M (10$\times$) & \$5M \\
Stage 2 & 1M & \$50M & \$500M (10$\times$) & \$50M \\
Stage 3 & 10M & \$500M & \$5B (10$\times$) & \$500M \\
Stage 4 & 100M & \$5B & \$50B (10$\times$) & \$5B \\
\midrule
\multicolumn{5}{l}{\textit{Base ARPU (\$100/yr)}} \\
Stage 1 & 100K & \$10M & \$150M (15$\times$) & \$15M \\
Stage 2 & 1M & \$100M & \$1.5B (15$\times$) & \$150M \\
Stage 3 & 10M & \$1B & \$15B (15$\times$) & \$1.5B \\
Stage 4 & 100M & \$10B & \$150B (15$\times$) & \textbf{\$15B} \\
\midrule
\multicolumn{5}{l}{\textit{Premium ARPU (\$200/yr)}} \\
Stage 1 & 100K & \$20M & \$400M (20$\times$) & \$40M \\
Stage 2 & 1M & \$200M & \$4B (20$\times$) & \$400M \\
Stage 3 & 10M & \$2B & \$40B (20$\times$) & \$4B \\
Stage 4 & 100M & \$20B & \$400B (20$\times$) & \textbf{\$40B} \\
\bottomrule
\end{tabular}
\end{small}
\end{center}

\textbf{Lux's 10\% equity, retained through the growth curve, is
worth between \$5M (Stage 1 Conservative) and \$40B (Stage 4
Premium).} Mid-case Base ARPU at 100M users implies \$15B of Lux
equity value.

% ─────────────────────────────────────────────────────────────────
\section*{VIII.\ Strategic Partnership Upside}

Lux brings substantial commercial value to a partnership-resolution
of this dispute. This section enumerates what the deal team can
offer as upside in any negotiation tier.

\subsection*{First Major Prospective Client: AvaTrade --- \$70B/month Trading Volume}

\textbf{Lux has direct access to AvaTrade as a prospective first
major client for the combined Lux + Counterparty venue.} The trading
flow contemplated by that engagement is on the order of \textbf{\$70
billion per month} routed through the joint architecture. AvaTrade
is an established multi-asset broker with operations across multiple
regulated jurisdictions; the framework underlying this prospective
arrangement is documented separately at
\texttt{\textasciitilde/work/lux/legal/partnerships/avatrade/Lux\_AvaTrade\_Strategic\_Partnership\_Agreement.pdf}.
\textit{The engagement is in advanced discussion under that framework
and is not yet a closed, signed customer commitment; the access to
AvaTrade is the Lux contribution, the venue and clearing stack is
the Counterparty contribution, and the joint take-rate is captured
together.}

\textbf{Fee economics.} The Counterparty's existing cost stack on
the securities-routing side is approximately \textbf{8 basis points
blended}, which is consumed internally between the Counterparty's
venue layer and the Alpaca clearing/custody relationship (split
approximately 4 bps / 4 bps, subject to volume-tier adjustment).
\textbf{Lux's take-rate sits on top of that internal cost stack and
is structurally flexible in the 25--50 bps range:}

\begin{itemize}[leftmargin=2em,itemsep=0.1em]
\item \textbf{Conservative} (25 bps Lux take) on \$70B/mo
$=$ \$175M/mo $=$ \textbf{\$2.1B/yr} in fee revenue to Lux from this
single relationship.
\item \textbf{Midpoint} (37.5 bps Lux take) on \$70B/mo $=$ \$262.5M/mo
$=$ \textbf{\$3.15B/yr}.
\item \textbf{Aggressive} (50 bps Lux take) on \$70B/mo $=$ \$350M/mo
$=$ \textbf{\$4.2B/yr}.
\item Even a \textbf{5 bps ``partnership uplift'' share-back} above
existing cost to each of Alpaca and the Counterparty (10 bps combined
incremental) is materially accretive to their per-trade economics and
preserves the bulk of the Lux take.
\end{itemize}

Customer-facing total fee at the conservative end (\textasciitilde 33
bps) remains competitive with multi-asset broker venues globally. At
Stage-2-plus volume routing, the AvaTrade contribution \textbf{exceeds
the Counterparty's current consolidated revenue base by approximately
an order of magnitude} and is, by itself, the strongest commercial
argument for partnership over litigation.

\subsection*{Crypto-Exchange Conversion at 8 Basis Points (Floor)}

A specific deliverable Lux brings under the partnership: \textbf{the
ability to offer crypto-exchange conversion at 8 basis points} ---
materially undercutting the incumbent fee structure at major US
crypto venues (Coinbase / Kraken / Gemini at 50--200 bps; Robinhood
PFOF-equivalent at 35--85 bps; institutional desks at 5--15 bps
floor but with size minimums). The 8-bps floor is enabled by Lux's
direct exchange and liquidity infrastructure and is a competitive
moat for the joint Lux + Counterparty venue.

\subsection*{OnyxPlus Access from Counterparty at Cost (Reciprocal Term)}

The Counterparty's OnyxPlus identity-verification product (a Telos JV in which the Counterparty holds the operating relationship) provides identity verification
infrastructure operating at the Lux ecosystem's internal cost basis
--- materially below merchant pricing for equivalent
KYC/AML/document/biometric services. Under the SCLA, \textbf{Lux
receives access to the Counterparty's OnyxPlus identity-verification
product at no charge or at the Counterparty's documented direct cost
(no markup, no commercial margin), not at market rates}. The annual savings depend on
the Counterparty's user-onboarding velocity but at Stage-2-plus
scale run into the tens of millions per year.

\subsection*{Co-Branding for Lux-Network Chains}

\textbf{All Lux-network L1/L2/L3/L4 chains receive explicit
co-branding on the Counterparty's customer-facing surfaces}: the
Counterparty's marketing, the venue's chain selector, the
SecurityToken issuance pages, and any institutional partnership
collateral. This brand co-positioning is a non-negotiable term
under the SCLA (\S XI).

\subsection*{Real-World Asset Pipeline --- Lux Brings Issuers; Counterparty Earns Listing + Trading Fees}

Lux operates the infrastructure layer behind multi-billion-dollar
RWA tokenization pipelines, and \textbf{brings the issuers}. The
pipeline is not theoretical --- Lux already operates a
\textbf{regulated Luxembourg digital-securities platform in
partnership with Creatrust (the licensed Luxembourg fund-manager
counterparty)}, on the same ERC-3643 / T-REX smart-contract suite
(\S IV.A(3) 2020--2023 milestone) that underlies the Counterparty's
ATS contract layer. The Luxembourg / Creatrust footprint covers
\textbf{tokenized commodities (uranium, gold, silver), tokenized
fund-of-funds (FoF) structures, and tokenized fund LP interests},
and provides both the CSSF-supervised regulatory wrapper and the
multi-year operational record that anchors Lux's issuer-side
credibility.

Under the SCLA, the Counterparty is the natural US-regulated
listing destination for that flow and \textbf{captures the
listing fees and the recurring trading-fee stream} on every asset
that goes up. The economic shape of the relationship is: Lux $=$
issuer origination, structuring, contract layer (ERC-3643 /
T-REX), and token-side cryptographic and chain-side machinery;
Counterparty $=$ regulated ATS / BD / TA listing venue + Alpaca
clearing + customer-facing distribution; \textbf{fees are split
per the SCLA fee schedule} (Lux take + Counterparty take + Alpaca
clearing cost). Asset classes in the pipeline:

\begin{itemize}[leftmargin=1.5em,itemsep=0.1em]
\item Tokenized US private-company equity and securities.
\item Tokenized real estate / NNN-SPV-wrapped property exposures.
\item \textbf{Tokenized commodities --- uranium, gold, silver, FX,
energy} (in production today on the Creatrust / Luxembourg
platform; SCLA migrates / cross-lists into the Counterparty's
US-regulated venue).
\item \textbf{Tokenized fund LP interests and fund-of-funds (FoF)
structures} (in production today on the Creatrust / Luxembourg
platform).
\item Tokenized fixed-income.
\item Tokenized music + IP royalty streams.
\item Tokenized art and collectibles ($>$\$1M lots).
\item \textbf{Substantial equity portfolio under tokenization
mandate (up to \$1 trillion in nominal value)} spanning AI
companies, resource / mining / energy companies, and other
sector-diverse portfolio holdings being prepared for tokenized
listing on the joint venue. \textbf{Every asset listed on the
Counterparty's venue generates (i) a one-time listing fee and (ii) a
recurring trading-fee stream over the asset's lifetime on-venue;}
the size and composition of the pipeline therefore materially shifts
the partnership's forward enterprise-value trajectory and is a
primary driver of the equity-share floor in \S X.
\end{itemize}

\textbf{Conservative estimate of incremental deal flow under the
SCLA: \$5--15 billion in nominal asset value over the first 12--24
months post-close on the multi-asset RWA pipeline above, plus the
phased onboarding of the \$1 trillion equity-portfolio mandate over
the medium term (24--60 months) as listings clear regulatory and
custodial preconditions. Recurring trading-fee revenue on the
seasoned book scales linearly with on-venue turnover.}

\subsection*{Reciprocal License to Counterparty ATS / BD / TA / Listed Securities --- Lux's Forward Venue Strategy}

\textbf{Lux Industries Inc.\ does not presently operate an
alternative trading system or broker-dealer in the United States.}
However, Lux's historical regulatory posture and its forward strategy
include venue licenses in multiple jurisdictions, and the reciprocal
license to the Counterparty's ATS / BD / TA / listed-securities
surfaces described in \S XI(9) is intended to give Lux operational
optionality across that footprint as it matures.

\textbf{Historical native support:} Lux (and Lux Partners Limited
as predecessor) previously operated regulatory arrangements in the
\textbf{Isle of Man (IOM)} (IOM Financial Services Authority
supervision) and \textbf{Luxembourg} (Commission de Surveillance du
Secteur Financier engagement). The IOM and Luxembourg footprints
provided native EU-passportable and offshore-fund-eligible regulatory
reach. Within the IOM footprint, the Lux family previously held,
through \textbf{CDAX Limited} (an IOM company in the Lux-affiliated
corporate group), an \textbf{IOM Class 8 money-transmission
licence}\footnote{IOM Financial Services Authority licensing class
covering money transmission and ancillary payment services; see
IOMFSA Regulated Activities Order.} --- an operational licensing
posture that establishes both the technical-operations and the
compliance-program track record relevant to the forward US (Reg
ATS / BD) and UK (FCA) applications and that is independently
diligenceable in the IOMFSA register. \textbf{Within the Luxembourg
footprint}, Lux operates a regulated digital-securities platform
in partnership with \textbf{Creatrust} (the licensed Luxembourg
fund-manager counterparty), running Lux's ERC-3643 / T-REX
contract suite to issue tokenized commodities (uranium, gold,
silver), tokenized fund-of-funds, and tokenized fund LP
interests --- the Luxembourg posture provides an independently
diligenceable CSSF-adjacent regulated venue and a multi-year
operational record for the same contract layer now embedded in
the Counterparty's ATS build.

\textbf{Prior BD + ATS development --- the 2018 ESX JV.}
Concurrent with the IOM Class-8 footprint, in 2018 \textbf{Hanzo,
Inc.\ pursued original BD + ATS licensure development through a
joint venture (the ``ESX JV'') with Ryan Kavanaugh and the Triller
group} (the social / short-form video group that subsequently
matured into the public Triller listing). The economic outcome of
the JV was that the Hanzo / Lux side received a \textbf{20\% stake
in ``Entertainment Stock X''} (the JV vehicle), and the BD + ATS
development work product generated within that JV was
\textbf{assigned / licensed forward into the Lux line of vehicles}
through the same corporate-identity chain documented at \S IV.A(1).
The ESX JV work product is therefore part of the upstream record
that materially \textbf{predates} the Counterparty's 2025--2026
build of its own ATS, supports the originality and
chain-of-attribution claims at \S IV.A(4) (patents) and Counts
I--III, and forms part of the Hanzo, Inc.\ source-licensor layer
described at \S IV.A(1B) --- which, again, was \textbf{never
licensed to the Counterparty}.

\textbf{Forward expansion:} Lux is actively evaluating venue-license
applications across the United States (Reg ATS), United Kingdom
(FCA MTF or equivalent), European Union (MiCA-compliant Trading
Platform), Singapore (MAS Capital Markets Services Licence +
Recognised Market Operator), and the UAE (VARA). Each license,
when obtained, becomes a venue at which Lux may operate the
Counterparty-licensed ATS / BD / TA technology under \S XI(9).

\textbf{Cross-listing rights:} Every security that lists on the
Counterparty venue is available for cross-listing on any Lux-network
or Lux-affiliate venue, and every security that lists on a Lux
venue is available for cross-listing on the Counterparty's venue.

\subsection*{Banking Licenses --- SF Private Bank Partnership and Global Coverage}

\textbf{Lux holds a strategic partnership with SF Private Bank
(``SFPB'') and, through SFPB and its corporate parent
\textbf{Supreme Fintech}, gains access to the consolidated SFPB /
Supreme Fintech licensed footprint across the populated continents.}
Adjacent partnerships with additional regulated financial
institutions fill remaining geographies. The combined licensed
coverage:

\begin{itemize}[leftmargin=1.5em,itemsep=0.1em]
\item \textbf{North America}: United States (federal banking partner,
FDIC-insured rails); Canada (chartered banking partner, CDIC-insured).
\item \textbf{Europe}: Sweden (Finansinspektionen-supervised, EU
passport); historical IOM (Financial Services Authority); historical
Luxembourg (CSSF); ongoing UK (FCA) and Switzerland (FINMA) expansion.
\item \textbf{Asia-Pacific}: Singapore (MAS Capital Markets Services
and banking); New Zealand (RBNZ-licensed deposit-taker); ongoing
Hong Kong (SFC type-1 / type-7) and Japan (JVCEA) expansion.
\item \textbf{Middle East}: UAE (VARA-supervised; DIFC / ADGM
counterparty relationships).
\item \textbf{Latin America}: regional partnerships under bilateral
arrangements; specific licensing on a case-by-case basis.
\item \textbf{Africa}: pilot arrangements; specific licensing on a
case-by-case basis.
\end{itemize}

Services Lux can deliver under the SCLA include Banking-as-a-Service
(BaaS), Payment Service Provider (PSP) rails, and \textbf{a custody
solution that materially undercuts the Fireblocks and BitGo
incumbents on price} (built on the same MPC infrastructure already
in the Counterparty's production).

\textbf{Estimated independent cost for the Counterparty to acquire
equivalent global license suite: \$57--140M in direct regulatory
and capital expenditure, plus 5--7 years of timeline.} See Schedule
C.

\subsection*{International Distribution: M7 and 1B+ End-User Reach}

Through Hanzo Inc., the broader corporate family, and the Lux
ecosystem partnerships, the M7 global initiatives provide
distribution access to more than 1 billion end-users across consumer,
institutional, and retail channels. Post-SCLA, Lux's international
counterparty network operates as a Reg-S-compatible distribution
channel for Counterparty-listed assets.

\subsection*{Hanzo Inc.\ Portfolio Precedents}

Hanzo Inc., the venture studio under which the Lux relationship has
historically operated, has applied the IP-in-kind founder-equity
model successfully to multiple spawned unicorn-class companies:

\begin{itemize}[leftmargin=1.5em,itemsep=0.1em]
\item \textbf{Damon Motorcycles} --- electric / smart-motorcycle company.
\item \textbf{Triller} --- short-form video / social platform.
\item \textbf{Bellabeat} --- women's health wearables.
\item \textbf{Unicorn Gold} --- precious-metals tokenization platform.
\item \textbf{Lux Network} --- itself a Hanzo-spawn outcome.
\item \textbf{Zoo Labs Foundation / Zoo Network} --- decentralized AI
research network.
\end{itemize}

Retained-parent-equity range across the portfolio: 10\%--40\%. The
10\% ask in the Counterparty is at the bottom of that internal range.

\subsection*{Private IP \& Patent Portfolio Access --- \$1B+ Aggregate Valuation}

In addition to the public-OSS-tier Lux modules and the LEL / LRL--PR
source-available tiers (\S IV.B), the Lux + Hanzo combined estate
includes a \textbf{private, closed-source, patent-bearing IP
portfolio} maintained outside the public repositories and not
present in the \texttt{luxfi/*} or \texttt{hanzoai/*} GitHub
organizations. Sources and scope (catalog being finalised in
\texttt{liquidity/review/patent\_catalog\_2026\_05\_25.md} adjacent
to this memorandum):

\begin{itemize}[leftmargin=1.5em,itemsep=0.1em]
\item \textbf{\texttt{\textasciitilde/work/luxcpp/}} --- closed-source
GPU / C++ acceleration tree: GPU-accelerated DEX matching engine
kernels; GPU-EVM execution kernels; GPU-accelerated cryptographic
primitives (BLS multi-pairing, KZG, FFT, MSM); multi-arch (CUDA /
Metal / ROCm / oneAPI) shim layer; AI-optimised bytecode
execution. Patent-bearing on multiple claims.
\item \textbf{Lux paper portfolio} (\texttt{\textasciitilde/work/lux/papers/})
--- the LP-* whitepaper series and topical specifications underlying
the patent filings enumerated at \S IV.A(4) (Lux side), including
\texttt{lp-020-quasar-consensus}, \texttt{lp-073-pulsar},
\texttt{lp-137-gpu-crypto-stack}, \texttt{lp-305} through
\texttt{lp-309} quantum-stack family, the \texttt{lux-ats-
architecture}, and the \texttt{gpu-evm-whitepaper}. Each paper
constitutes a defensive-publication or filable-patent record.
\item \textbf{Hanzo paper portfolio} (\texttt{\textasciitilde/work/hanzo/papers/})
--- HIP-* and thematic papers including \texttt{hanzo-aci},
\texttt{hanzo-agent-grpo} / \texttt{hanzo-agent-sdk},
\texttt{hanzo-agi-platform}, \texttt{hanzo-ai-chain},
\texttt{hanzo-consensus-ai}, \texttt{hanzo-cloud-
infrastructure}, \texttt{hanzo-blockchain-platform}, the
\texttt{jin/} multimodal subtree, and the \texttt{defense/}
subtree (assured-networking, cross-domain, countermeasures,
edge-AI/ML). Each paper constitutes a defensive-publication or
filable-patent record on the Hanzo side.
\item \textbf{Adjacent private repositories} across
\texttt{\textasciitilde/work/} bringing pre-existing trading-
platform, exchange, AI / ML, and infrastructure work product
into scope: \texttt{\textasciitilde/work/beyondequity},
\texttt{\textasciitilde/work/beyondalpha}, and other Lux / Hanzo
internal monorepos. (Counsel: complete enumeration during
diligence.)
\end{itemize}

\textbf{Aggregate valuation: in excess of \$1 billion.} The
combined Lux + Hanzo private-IP and patents estate, valued on a
comparable-transaction basis against precedent crypto and AI
infrastructure IP transactions (and against the Counterparty's
own current \$1.1B raise valuation, which materially reflects the
contribution of this IP via Mr.\ Kelling's Q4-2025 rebuild), sits
\textbf{north of \$1 billion in aggregate value} before any
consideration of forward licensing or commercialisation revenue.
That valuation is a function of (i) the number and breadth of
patent-bearing inventions; (ii) the demonstrated reusability of
the OSS-public layer to extract enterprise value (the
Counterparty's own raise being the immediate evidence); and (iii)
the rarity of an integrated cryptographic + post-quantum + GPU +
AI + exchange-engine portfolio at this depth.

\textbf{SCLA grant scope.} Under the SCLA, the Counterparty
receives access to this private portfolio \textbf{at the Maximum
(A+) tier only}, and at progressively more limited subsets at
lower tiers (see \S X feature-unlock table). At any tier, the
SCLA grant runs to the Counterparty as licensee, subject to (a)
non-sublicensability outside the Counterparty's regulated venues,
(b) the standard LEL / LRL--PR-style ``Authorized Network'' /
``Descending Chain'' criterion (\S IV.B), and (c) the SCLA's
material-breach termination provisions. \textbf{The Hanzo, Inc.\
co-licensor consent (\S XI ¶20, ¶22A) is required for the
Counterparty to access the Hanzo-side portion of the portfolio.}

\subsection*{The \$100B--\$1T Target Institution}

The combined Lux + Counterparty vision is the construction of a
\$100-billion-to-\$1-trillion-class institution at the intersection
of tokenized real-world assets, post-quantum-secure infrastructure,
regulated US securities markets, and global distribution. The deal
team should treat the negotiation not as a zero-sum settlement but
as the founding-document of a partnership creating value materially
in excess of the negotiating range.

% ─────────────────────────────────────────────────────────────────
\section*{IX.\ Partnership Now vs.\ Delay Analysis}

\textbf{Forces favoring delay:} equity appreciation; royalty accrual;
current-raise disgorgement expansion; pattern-of-conduct evidence.

\textbf{Forces favoring settlement now:} statutory damages caps don't
grow with delay (DTSA at 2$\times$; patent at 3$\times$); dilution risk
on Lux's equity component; operational risk (engineer-around); evidence
preservation; DTSA 3-year limitations from discovery; Lux operating
runway; negotiating leverage compounds short-term then evaporates if
the Counterparty hits trouble; public visibility window.

\textbf{Recommendation:} \textbf{Partner now, with strong protective
terms} --- significant cash component, 10\%+ equity with weighted-
average broad-based anti-dilution, full ratchet on down-rounds, pro-
rata participation, tag-along and co-sale rights, 1+1 mandatory board
seats, 99-year exclusive USA-market license to the Counterparty
(conditional), patent grant, AvaTrade routing, 8 bps conversion floor,
OnyxPlus access from Counterparty at zero / documented cost, co-branding for Lux chains, and
reciprocal permanent license to Counterparty ATS / BD / TA / listed-
securities / cross-listings (\S XI). \textbf{Statutory damages caps
do not grow with delay; the equity does.} Locking equity now with
strong anti-dilution captures the future-growth narrative without
the operational risks of continued reliance on a deteriorating
informal arrangement.

% ─────────────────────────────────────────────────────────────────
\section*{X.\ Partnership Triangulation (Equity-Share + Cash + Governance + Reciprocal Commitments)}

\subsection*{Posture}

This is a \textbf{partnership-structuring matrix, not a litigation-
settlement matrix}. The deal team's job is to land a structure that
both sides can live with --- one that regularises the IP relationship,
preserves the Counterparty's operational continuity (including Mr.\
Nandy's team's role), gives the Counterparty's current-round investors
a clean cap table, and gives Lux Industries Inc.\ a substantive
ongoing economic interest. Litigation (\S V--VI) is the alternative
posture if a partnership cannot be reached; the available exposure
range is \$80M--\$890M cash and is documented as leverage, not as
the operative ask.

\subsection*{Partnership Sizing Tiers (10\%--50\% equity range)}

The deal range is \textbf{10\% equity (minimum substantive
partnership) to 50\% equity (full equal-partner recapitalisation)}.

\textbf{Cash and equity are inversely related to a constant total-
value target.} The reference point is the maximum-equity tier (50\%
$\times$ \$1.1B = \$550M, no cash required because equity alone
satisfies the target). At lower equity percentages, \textbf{cash
fills the gap}: cash $=$ (50\% $-$ agreed equity \%) $\times$ \$1.1B.
\textbf{Equity $\uparrow$ $\Rightarrow$ cash $\downarrow$.} The
Counterparty chooses where in the range to land based on its
preference between existing-shareholder dilution (equity-heavy
preserves cash) and cash outflow (cash-heavy preserves cap-table).
Each tier additionally carries an escalating governance package and
broader reciprocal commercial commitments (right column).

\begin{center}
\begin{small}
\begin{tabular}{p{1.9cm} r r r p{4.8cm}}
\toprule
\textbf{Tier} & \textbf{Equity} & \textbf{Cash} & \textbf{Total nominal} & \textbf{Governance + reciprocal} \\
\midrule
\textbf{Maximum (M)}    & \textbf{50\%} & \textbf{\$0}     & \textbf{\$550M} (\$550M equity + \$0 cash) & Governance parity; equal-partner recap; full reciprocal license; international M7 distribution; closed-source \texttt{luxcpp} GPU access \\
\textbf{Stretch (S)}    & \textbf{35\%} & \textbf{\$165M}  & \textbf{\$550M} (\$385M + \$165M) & Board seats (2); above + co-branding for Lux-network chains; AvaTrade flow routing; 8\,bps crypto-conversion floor \\
\textbf{Recommended (R)}& \textbf{25\%} & \textbf{\$275M}  & \textbf{\$550M} (\$275M + \$275M) & Board seat (1); above + BaaS / SF Private Pay access; joint GTM commitments \\
\textbf{Floor (F)}      & \textbf{10\%} & \textbf{\$440M}  & \textbf{\$550M} (\$110M + \$440M) & Board observer; LEL/LRL-PR perpetual license to Counterparty; reciprocal BD/ATS/TA access to Lux \\
\bottomrule
\end{tabular}
\end{small}
\end{center}

\textbf{The math:} cash $=$ (50\% $-$ agreed equity \%) $\times$
\$1.1B. So 50\% equity $\Rightarrow$ \$0 cash; 35\% $\Rightarrow$
\$165M; 25\% $\Rightarrow$ \$275M; 10\% $\Rightarrow$ \$440M. Every
tier delivers the same \textbf{\$550M total nominal value} to Lux,
which is the floor of Lux's claim (50\% of enterprise value derived
from Lux IP).

\textbf{Practical constraint --- cash sequencing.} The Counterparty's
post-current-raise cash position will not support large lump-sum cash
payments at the lower-equity tiers in a single payment. The deal team
should sequence the cash leg as a \textbf{multi-year schedule}
(e.g.\ 30\% at close, balance over 24--36 months from operating cash
flow, secured by a security interest in Lux's pledged equity stake or
a default-to-additional-equity conversion clause). A multi-year cash
schedule shifts the practical landing zone toward the higher-equity
tiers --- which is the natural commercial outcome and which is what
Lux wants (more equity participation in long-term value).

\subsection*{Pricing Comp --- Anchor on the Current Round}

The Counterparty's current \$75M raise is being subscribed at \$1.1B
valuation (Forum Markets lead, ~\$37.5M soft-circled). Equity issued
to Lux at the same valuation is priced on the same per-share basis
as the new-money investors. This is the clean precedent transaction
to anchor on --- no separate valuation exercise required, no
contested damages multiplier, no dispute over the basis of the
equity value.

\subsection*{Feature-Unlock by Settlement Tier}

\begin{center}
\begin{small}
\begin{tabular}{p{6cm} ccccc}
\toprule
\textbf{Forward-value feature} & \textbf{S--} & \textbf{S} & \textbf{M} & \textbf{B} & \textbf{A+} \\
\midrule
LEL/LRL--PR perpetual license & $\checkmark$ & $\checkmark$ & $\checkmark$ & $\checkmark$ & $\checkmark$ \\
Patent grant on covered IP & $\checkmark$ & $\checkmark$ & $\checkmark$ & $\checkmark$ & $\checkmark$ \\
Maintenance \& support of upstream OSS & & $\checkmark$ & $\checkmark$ & $\checkmark$ & $\checkmark$ \\
99-year exclusive USA-market license & & $\checkmark$ & $\checkmark$ & $\checkmark$ & $\checkmark$ \\
1 Lux + 1 independent director seat & & $\checkmark$ & $\checkmark$ & $\checkmark$ & $\checkmark$ \\
RWA pipeline routing into Counterparty & & $\checkmark$ & $\checkmark$ & $\checkmark$ & $\checkmark$ \\
AvaTrade flow routing into venue & & $\checkmark$ & $\checkmark$ & $\checkmark$ & $\checkmark$ \\
8 bps crypto-conversion floor & & $\checkmark$ & $\checkmark$ & $\checkmark$ & $\checkmark$ \\
OnyxPlus access from Counterparty at zero / documented cost & & $\checkmark$ & $\checkmark$ & $\checkmark$ & $\checkmark$ \\
Co-branding for Lux-network chains & & $\checkmark$ & $\checkmark$ & $\checkmark$ & $\checkmark$ \\
Reciprocal license to ATS/BD/TA & & & $\checkmark$ & $\checkmark$ & $\checkmark$ \\
SFPB banking-license access (BaaS, PSP) & & & $\checkmark$ & $\checkmark$ & $\checkmark$ \\
Custody undercutting Fireblocks & & & $\checkmark$ & $\checkmark$ & $\checkmark$ \\
International M7 distribution channels & & & & $\checkmark$ & $\checkmark$ \\
Private Lux portfolio (\texttt{luxcpp} GPU + LP-* papers) & & & & $\checkmark$ & $\checkmark$ \\
Private Hanzo portfolio (\texttt{jin}, \texttt{defense}, ACI, agent SDK papers) & & & & & $\checkmark$ \\
Full \texttt{\textasciitilde/work/*} private estate (\$1B+ valuation) & & & & & $\checkmark$ \\
Joint go-to-market programs & & & & $\checkmark$ & $\checkmark$ \\
\bottomrule
\end{tabular}
\end{small}
\end{center}

\subsection*{Partnership vs.\ Litigation Triangulation}

\begin{center}
\begin{small}
\begin{tabular}{lrr}
\toprule
\textbf{Path} & \textbf{Lux receives} & \textbf{Counterparty pays} \\
\midrule
\textbf{Partnership Maximum (M, 50\%)}        & \textbf{\$550M equity + \$0 cash + governance parity + full reciprocal} & \textbf{\$550M nominal} \\
\textbf{Partnership Stretch (S, 35\%)}        & \textbf{\$385M equity + \$165M cash + 2 board seats + GTM + co-branding} & \textbf{\$550M nominal} \\
\textbf{Partnership Recommended (R, 25\%)}    & \textbf{\$275M equity + \$275M cash + board seat + BaaS access} & \textbf{\$550M nominal} \\
\textbf{Partnership Floor (F, 10\%)}          & \textbf{\$110M equity + \$440M cash (multi-year) + board observer + LEL grant} & \textbf{\$550M nominal} \\
\midrule
\multicolumn{3}{l}{\textit{Alternative path (litigation; documented as leverage, not as the operative ask):}} \\
Litigate: Small (Path A)                  & \$80M cash & \$80M cash + no relationship \\
Litigate: Medium (Path B mixed)           & \$215M cash & \$215M cash + no relationship \\
Litigate: Big (Path B full)               & \$590M cash & \$590M cash + no relationship \\
Litigate + Opp.\ Cost (Big)               & \$690--890M cash & idem + no relationship + reputational \\
\bottomrule
\end{tabular}
\end{small}
\end{center}

\textbf{Every partnership tier strictly dominates litigation} on a
forward-value basis because (i) the equity component participates in
the Counterparty's upside, (ii) the reciprocal commercial commitments
generate operating revenue for both sides, (iii) Mr.\ Nandy's team
remains in the build, and (iv) Lux's IP is regularised cleanly under
the LEL/LRL-PR perpetual license that ships with every tier. The
litigation rows are documented for negotiating purposes and as the
alternative posture if partnership terms cannot be agreed; they are
not the operative ask.

% ─────────────────────────────────────────────────────────────────
\section*{XI.\ Mandatory Terms (Lux's Floor in Any Partnership)}

The following terms are non-negotiable. Any partnership outcome at
any of the four tiers (Floor / Recommended / Stretch / Maximum) must
include each of them.

\textbf{0.} \textbf{Key-person retention --- Mr.\ Zach Kelling
remains as Counterparty Chief Technology Officer.} Lux's equity
stake is rendered substantially worthless without Mr.\ Kelling's
continued technical leadership of the Counterparty platform; the
current Counterparty operational team absent Mr.\ Kelling is not
viewed as competent to maintain or extend the rebuilt stack that
underpins the Counterparty's enterprise value (see \S IV.E ¶25 and
\S IV.F). The partnership therefore requires:
\begin{itemize}[leftmargin=2em,itemsep=0.1em]
\item Mr.\ Kelling's continued tenure as Counterparty CTO for a
minimum term of \textbf{four (4) years} from closing, with defined
good-leaver / bad-leaver triggers and a 12-month garden-leave / non-
compete on the Counterparty side if termination is without cause.
\item Mutual general releases of Mr.\ Kelling individually from
Counts I, II, III, IV, VI, VII, IX, XI, and XII of this memorandum,
conditional on full execution of the partnership agreement and
ongoing performance under it.
\item Carve-out of Mr.\ Kelling from any litigation Lux may bring
against other Counterparty officers (Choi / Trombley / Church /
Sonsurkar) if partnership terms are achieved.
\item Compensation structure (cash + Counterparty equity grants)
benchmarked against the comparable CTO market and adjusted to align
his continuing economic interest with both Lux (via partnership
equity) and the Counterparty (via direct compensation).
\item Bad-leaver clawback triggers tied to material breach of
post-settlement obligations.
\end{itemize}

\textbf{1.} \textbf{Board composition --- expansion to a minimum of
seven (7) directors total, with structured composition.} The
Counterparty's existing three-director board failed to identify the
licensing mismatch over the relevant period despite having a
maintainer-access-empowered CSO whose duty included precisely that
review. \textbf{The 3-person board structure is, in light of that
record, an inadequate governance vehicle going forward}, and the
disinterested-investor case for board expansion is unambiguous.

The mandatory composition is:

\begin{itemize}[leftmargin=2em,itemsep=0.1em]
\item \textbf{Three (3) existing director seats retained} from the
current board.
\item \textbf{One (1) Lux director seat}, for Lux Industries Inc.'s
direct representative (Mr.\ Pahlavi's designate).
\item \textbf{Three (3) mutually-agreed independent director seats},
each nominee with no prior commercial or personal relationship to
either Lux or the Counterparty, serving as the ongoing governance
anchor of the SCLA.
\end{itemize}

Effective at closing. \textbf{No ownership-maintenance condition
attaches to the Lux seat.} The structure ensures (a) no single
faction (existing officers, Lux, or independents) holds majority
voting power; (b) the independent block of three directors provides
the disinterested-quorum threshold for any related-party transaction
under Delaware General Corporation Law \S 144; and (c) the
Counterparty's institutional governance moves from informal,
founder-controlled posture to formal, multi-faction supervision
appropriate to a venue valued in the \$1B+ range with FINRA / SEC
exam exposure.

In the alternative (a 7-director soft floor that Lux will accept
under appropriate compensating concessions), the composition may
shift to: 2 existing seats retained + 2 Lux directors + 3 independent
seats. This variant displaces one of the existing seats and
strengthens Lux's voting position; Lux's deal team may invoke this
alternative as leverage if the Counterparty resists the principal
structure above.

\textbf{2.} \textbf{99-year exclusive USA-market license to the
Counterparty} for the LEL/LRL--PR portfolio in scope. Conditional on
maintenance of all SCLA covenants. Material breach terminates
exclusivity. Outside the USA, Lux retains the right to license to
other parties.

\textbf{3.} \textbf{Anti-dilution: weighted-average broad-based,
with full ratchet on any down-round below the Lux issuance price.}

\textbf{4.} \textbf{Pro-rata participation right} on all future
Counterparty preferred financings.

\textbf{5.} \textbf{Tag-along and co-sale rights} on any founder
transfer exceeding 1\% of fully-diluted equity.

\textbf{6.} \textbf{Information rights}: quarterly and annual
financials; all board materials and minutes; all share-issuance and
option-grant resolutions; counsel and accountant access under NDA.

\textbf{7.} \textbf{Patent grant} on all current and future Lux
patents covering the licensed IP. Worldwide; perpetual; royalty-free
under the SCLA; expressly inclusive of patents not yet filed at SCLA
execution.

\textbf{8.} \textbf{Maintenance and support of upstream Lux OSS}:
continuing on the shared-engineering basis as today, no separate
maintenance fee.

\textbf{9.} \textbf{Reciprocity --- Permanent license to Counterparty
ATS, BD, TA, listed securities, and cross-listings.} As symmetric
consideration for Lux's perpetual IP license, the Counterparty grants
Lux Industries Inc.\ (and any Lux-controlled affiliate or Lux-network
L1/L2/L3/L4 chain maintaining Lux consensus and Warp interop):

\begin{itemize}[leftmargin=2em,itemsep=0.1em]
\item \textbf{Permanent, royalty-free, non-exclusive license to the
Counterparty's ATS technology and operational know-how}, sufficient
for Lux to operate (or sublicense to a Lux-controlled regulated
entity) an ATS in any jurisdiction where Lux holds or obtains the
necessary venue license.
\item \textbf{Permanent, royalty-free, non-exclusive license to the
Counterparty's BD technology and know-how} on equivalent terms.
\item \textbf{Permanent, royalty-free, non-exclusive license to the
Counterparty's TA technology and know-how} on equivalent terms.
\item \textbf{Reciprocal listing and cross-listing rights}: every
security listed on the Counterparty venue is available for
cross-listing on any Lux venue, and every security listed on a Lux
venue is available for cross-listing on the Counterparty venue, on
no-less-favorable terms.
\item \textbf{Sublicensing rights}: Lux may sublicense to Lux-
controlled jurisdictions and operating affiliates, including the
IOM, Luxembourg, FCA, MAS, VARA, MiCA, and FINMA jurisdictional
expansions described in \S VIII.
\end{itemize}

The reciprocal grant runs 99 years, conditional on continued SCLA
performance.

\textbf{10.} \textbf{Co-branding for Lux-network chains.} All Lux-
network L1/L2/L3/L4 chains receive explicit co-branding on the
Counterparty's customer-facing surfaces (marketing, venue chain
selector, SecurityToken issuance pages, and institutional partnership
collateral).

\textbf{11.} \textbf{OnyxPlus access from Counterparty at zero or
documented cost.} Lux receives access to the Counterparty's OnyxPlus
KYC / AML / document / biometric identity-verification product
(operated by the Counterparty as a Telos JV) at no charge or at the
Counterparty's documented direct cost (no markup, no commercial
margin). This is a reciprocal commitment from the Counterparty to
Lux, not the other direction.

\textbf{12.} \textbf{AvaTrade routing}: the AvaTrade flow on which
Lux holds the direct access (under the framework documented in
\S VIII) is routed through the joint venue post-SCLA; the
Counterparty receives the proportional benefit of the trading
volume on the joint-take terms (25--50 bps Lux take, 8 bps
internal Counterparty / Alpaca cost stack; see \S VIII First Major
Prospective Client). Closing of the AvaTrade customer commitment
is a separate workstream and not a condition precedent to the
SCLA.

\textbf{13.} \textbf{8 bps crypto-conversion floor}: Lux's
infrastructure enables the joint venue to offer crypto-exchange
conversion at 8 basis points (floor); the Counterparty's existing
crypto-conversion pricing is repriced to this floor under the SCLA.

\textbf{14.} \textbf{Engineering commitment}: the Counterparty
activates the Lux Warp precompile (canonical address
\texttt{0x0200000000000000000000000000000000000005}) on every Liquid
chain at the next chain-upgrade window. New Liquid chains include
the \texttt{warpConfig} block in genesis.

\textbf{15.} \textbf{Litigation hold and document preservation}: the
Counterparty agrees to preserve all records, communications, and
code-review artifacts relating to the licensing posture from 2014 to
the SCLA execution date.

\textbf{16.} \textbf{Independent investigation cooperation}: the
Counterparty cooperates fully with any independent investigation of
officer conduct undertaken by the disinterested directors or their
appointed counsel.

\textbf{17.} \textbf{Whistleblower / non-retaliation protection ---
Mr.\ Tridib Nandy (Counterparty VPE; Lux CTO).} The Counterparty
agrees that Mr.\ Nandy, who identified and elevated this matter to
the Lux Industries Inc.\ board through his dual-role
\texttt{luxfi/*}-admin visibility and concurrent VPE seat, shall
suffer no retaliation, demotion, compensation reduction, or
constructive termination from the Counterparty side as a result of
his role in this matter. Specifically:
\begin{itemize}[leftmargin=2em,itemsep=0.1em]
\item Mr.\ Nandy's VPE role and reporting line at the Counterparty
shall be preserved on no-less-favourable terms than as of the day
preceding the elevation of this matter.
\item Indemnification from the Counterparty for any claim asserted
against Mr.\ Nandy personally arising from his role in identifying
or escalating the matter, to the maximum extent permitted by
applicable corporate law.
\item Mutual general release of Mr.\ Nandy individually from any
counterclaim the Counterparty might assert (conditional on full
execution of the partnership agreement and ongoing performance).
\item Carve-out of Mr.\ Nandy from any litigation the Counterparty
may bring against Lux or any other party arising from this matter.
\end{itemize}

\textbf{18.} \textbf{Lux engineering team non-solicit / non-
retaliation.} The Counterparty agrees, for a period of 36 months
post-SCLA execution, not to solicit, recruit, or otherwise induce
the departure of any member of the Lux engineering team identified
in this matter (a closed list of named individuals to be appended to
the SCLA), and not to retaliate against any such individual for any
role in the discovery, escalation, or remediation of this matter.
General-market hiring through bona fide third-party recruiters or
in response to unsolicited applications by the individual is
excluded.

\textbf{19.} \textbf{D\&O indemnification of Lux's officers and
directors for raising this matter.} The Counterparty agrees, by
covenant in the SCLA, to release Lux Industries Inc.\ and its
present and former officers and directors (including but not
limited to Messrs.\ Pahlavi, Dupont, Donohue, Nandy, and Kelling)
from any direct or derivative claim arising from the Counterparty's
use of Lux IP through the SCLA execution date or arising from any
party's role in identifying, escalating, or seeking to remediate
the matter. The release is mutual and integral to the partnership
outcome; refusal of the release by the Counterparty is a deal-
breaker.

\textbf{20.} \textbf{Lux and Hanzo joint warranty of unencumbered
title to the licensed IP through the SCLA grant.} Lux Industries
Inc.\ warrants that, as of the SCLA execution date, the corporate-
identity chain \textbf{Monoid LLC (and earlier research vehicles,
2008+) $\rightarrow$ Hanzo, Inc.\ (2016) $\rightarrow$ Lux
Partners Limited (IOM) $\rightarrow$ Lux Industries Inc.}
carries unencumbered title to the LEL- and LRL-PR-licensed IP
being granted under the SCLA. \textbf{Hanzo, Inc.}, as the
source-licensor of the Hanzo-origin contribution layer
(\S IV.A(1B)), additionally grants under the SCLA a parallel
royalty-free permanent licence to the Hanzo-source portion of the
contributed work product, so that the SCLA regularises both the
Lux-licensor and the Hanzo-licensor gaps in a single instrument.
Reasonable and customary IP-warranty wording; assignment
instruments between successive corporate vehicles to be referenced
in the SCLA exhibits.

\textbf{21.} \textbf{Permanent access for Hanzo, Inc.\ and Zoo
(Zoo Labs Foundation) to the Counterparty's protocol and venue.}
The SCLA shall grant Hanzo, Inc.\ and the Zoo Labs Foundation
each a \textbf{permanent, royalty-free, non-exclusive access
right} to the Counterparty's protocol APIs, regulated venue
surfaces (ATS / BD / TA), and listed-securities cross-listing
mechanisms, for their respective AI / infrastructure (Hanzo) and
AI-research / decentralized-science (Zoo) purposes. This access
is in addition to Lux's primary equity / cash / reciprocal-licence
package and is not a substitute for any other term. The
permanent-access grant runs with the LEL grant and survives any
change of control or recapitalisation of the Counterparty.

\textbf{22.} \textbf{Hanzo, Inc.\ as named beneficiary of the
SCLA.} Hanzo, Inc.\ shall be named as an express third-party
beneficiary of the SCLA's reciprocal-licence, OSS-contribution
recognition, and permanent-access provisions, with its own
standing to enforce those provisions independent of Lux Industries
Inc.'s enforcement rights.

\textbf{22A.} \textbf{Hanzo, Inc.\ co-licensor consent and mutual
release.} Because the Hanzo-source IP contribution layer
(\S IV.A(1B), Counts I--III, V--VII, XII) is independently held
by Hanzo, Inc.\ and was never licensed to the Counterparty, the
SCLA shall include:

\begin{itemize}[leftmargin=2em,itemsep=0.1em]
\item \textbf{Hanzo, Inc.\ as a signing party} to the SCLA in its
own name, with parallel grant of a perpetual, royalty-free
licence to the Hanzo-source portion of the contributed work
product (matching the LEL / LRL--PR perpetual licence ¶2 grant
from Lux on the Lux-source portion).
\item \textbf{Hanzo, Inc.\ co-licensor consent to the Counterparty's
access} to the Hanzo-side portion of the private IP \& patents
portfolio described in \S VIII (jin, defense, ACI, agent SDK, AGI
platform, cloud / blockchain platform, consensus-AI papers, the
Hanzo Inc.\ patent portfolio at \S IV.A(4)).
\item \textbf{Mutual general release of Hanzo, Inc.\ and its
present and former officers, directors, employees, advisors, and
consultants} (including but not limited to Mr.\ Kelling in his
Hanzo, Inc.\ consulting capacity, see \S IV.A(1C)) from any
direct or derivative claim arising from the Counterparty's use of
Hanzo-source IP through the SCLA execution date.
\item \textbf{Lux + Hanzo joint release of the Counterparty} on
final execution and ongoing performance under the SCLA, in
respect of all Counts pleaded jointly or severally in this
memorandum. Refusal of the joint release by the Counterparty is
a deal-breaker.
\end{itemize}

The Hanzo, Inc.\ release on its own behalf, and the parallel
Hanzo, Inc.\ consent for its underlying IP to flow to the
Counterparty under the SCLA, are conditions precedent to closing.
Without them, the Hanzo-source layer remains independently
actionable against the Counterparty (\S IV.A(1B); Counts I--III,
V--VII, XII), and the partnership outcome cannot be a clean
final disposition.

% ─────────────────────────────────────────────────────────────────
\section*{XII.\ Special Note Re: Mr.\ Zach Kelling}

\textbf{Mr.\ Kelling occupies a position requiring particular care
in Lux's partnership and (alternative-path) litigation strategy.}
The factual context established in \S IV(E) and \S IV(F) bears
directly here:

\begin{itemize}[leftmargin=1.5em,itemsep=0.1em]
\item Mr.\ Kelling \textbf{departed Lux Industries Inc.}\ and is
\textbf{currently Chief Technology Officer of the Counterparty
(Liquidity.io) only} (not concurrently a Lux officer or director).
\item He was the \emph{remediator} of the pre-Kelling Counterparty
technical and security failures (the broken time-price priority
matching; the \$40,000 ACH-exploit theft) and the architect of the
rebuilt platform that supports the current \$75M / \$1.1B-valuation
raise, including personally introducing Alpaca as the clearing BD
and developing the Alpaca-integration exchange tech over Q4 2025
(\S IV.A(3)).
\item \textbf{AI-assisted-development mitigating context}: Mr.\
Kelling and his team conducted the rebuild substantially using
AI-assisted development tooling at the rapid pace characteristic
of late-2025 / early-2026 AI-augmented engineering. Whether the
under-recognition of the Lux LICENSE markers carried into the
imported modules was intentional or (as is more likely given the
AI-driven pace) unintentional is a question for fact development;
in either case, Mr.\ Kelling has indicated to Lux that he now
sides with the Lux team on the clean precedent that the Lux IP
required a commercial licence and that the appropriate disposition
is partnership regularisation, not adversarial litigation.
\item Mr.\ Kelling is not the author of the disclosure conduct
underlying Counts IX and XI; the duty to disclose the licensing
mismatch to investors sat with the Counterparty's CEO and
strategy/CSO chain (\S IV.D, \S V Count V).
\end{itemize}

\textbf{Litigation-strategy implications:}

\begin{itemize}[leftmargin=1.5em,itemsep=0.1em]
\item Mr.\ Kelling should \emph{not} be named as a defendant in any
Lux suit. The conduct alleged is not his.
\item Mr.\ Kelling is likely a material witness for Lux on (a) the
authorship and ownership of the Lux IP, (b) the absence of any
executed commercial license, and (c) the operational-displacement
narrative explaining why the mismatch was not raised internally.
\item Counsel should consider a litigation-cooperation agreement,
witness-protection arrangement, or formal indemnification of Mr.\
Kelling (as a non-defendant witness) to ensure his continued
availability without undue personal exposure.
\item In any public communication, Lux should take care not to
attribute the misconduct to Mr.\ Kelling. The narrative should be
specific to the Counterparty officers actually responsible.
\item In the SCLA itself, Lux should provide for Mr.\ Kelling's
protective treatment, including (a) a release of any claims Lux
might assert against him personally, (b) acknowledgment of his
historical authorship, and (c) a ROFR / similar mechanism if his
Counterparty role generates IP that should flow back to Lux upstream.
\end{itemize}

\textbf{Mr.\ Kelling can serve as a productive intermediary in the
Lux--Counterparty negotiation}, given his founder lineage with Lux,
his current role as Counterparty CTO, his stated alignment with the
Lux position on the need for licensing regularisation, and his
status as the indispensable key-person on the Counterparty's
operational side (\S XI ¶0). Mr.\ Dupont should consider engaging
Mr.\ Kelling directly in informational meetings during the
negotiation phase, subject to outside counsel's oversight, with
the explicit framing that key-person retention (\S XI ¶0) is on
the table as part of the partnership outcome.

% ─────────────────────────────────────────────────────────────────
\section*{XIII.\ Procedural Recommendations}

\subsection*{A.\ Recommended Approach to Counterparty (Partnership-First)}

The operative posture is \textbf{partnership-first / litigation-as-
alternative}. The sequencing the deal team should run:

\begin{enumerate}[leftmargin=1.5em,itemsep=0.1em]
\item \textbf{Quiet inbound to Counterparty CEO (Mr.\ Choi)} ---
chair-to-chair / counsel-to-counsel introduction of this matter
under a friendly framing, leading with the partnership offer
(§X Recommended tier as opener: 25\% equity + ~\$275M cash to
$\$$550M nominal). Frame the conversation as "we identified a
licensing mismatch that needs regularisation; here is how we
propose to regularise it." Do not lead with Counts or damage
ranges.
\item \textbf{Parallel quiet investigation track} --- continue
outside-counsel diligence on the on-disk evidence (commit history,
LICENSE-marker travel into the Counterparty tree, current-raise
disclosure materials, Sonsurkar / Trombley / Choi communication
record) so the leverage is preserved if the partnership track
stalls. Litigation hold notice (§XI ¶15) issued early.
\item \textbf{Triangulate to Recommended (R, 25\%)} --- this is the
landing zone Lux should anchor on. Floor (10\%) is acceptable only
with substantial non-equity package (BD/ATS/TA reciprocal license,
BaaS access, multi-year cash schedule) that NPV-equivalents the
equity gap. Stretch (35\%) and Maximum (50\%) are upside if the
Counterparty wants to preserve cash via additional equity.
\item \textbf{Multi-year cash schedule} --- the cash component at
the lower-equity tiers is large in nominal terms (\$275M at R,
\$440M at F) and not deliverable in a single payment from current
Counterparty liquidity. Sequence as 30\% at close, balance over
24--36 months from operating cash flow, secured by a security
interest in Lux's pledged equity stake or a default-to-additional-
equity conversion clause (cap-table dilution mechanically
backfills if cash schedule slips).
\item \textbf{Mutual releases / symmetric protections} ---
partnership outcome ships with: §XI ¶0 Kelling key-person
retention with mutual releases; §XI ¶17 Nandy whistleblower /
non-retaliation; §XI ¶18 Lux engineering team non-solicit; §XI
¶19 D\&O indemnification; §XI ¶20 Lux IP warranty. Refusal of
any of these is a deal-breaker.
\item \textbf{Reciprocal commercial commitments} --- the §X tier
governance / reciprocal package (board seats, LEL grant, BaaS,
joint GTM, co-branding, ATS/BD/TA reciprocal license) is
substantively load-bearing, not aesthetic. Each item is what gives
the deal forward operating value to Lux beyond the equity stake
itself.
\item \textbf{Escalation ladder if partnership stalls}: (i) reduce
the equity ask one tier (do NOT raise it) to surface whether
substance-of-objection is equity-size or something else; (ii) if
still stalled, ramp the quiet-investigation track to a sealed
filing (preserves leverage, doesn't go public yet); (iii) if no
movement within 90 days from first quiet inbound, file publicly
on the alternative-posture Counts.
\item \textbf{Communications discipline} --- no public statements
during partnership-first period. Mr.\ Pahlavi or Mr.\ Donohue (GC)
clears every external communication. Internal-comms reminder
issued to Lux engineering team about confidentiality of pending
matter.
\end{enumerate}

\textbf{Time-to-deal target}: 60--90 days from first quiet inbound
to executed partnership term sheet, conditional on Counterparty
engagement. The current Counterparty raise close is a natural
forcing function and should be used as the deal-window anchor.

\subsection*{B.\ Decision Chain}

\textbf{Decision chain:}

\begin{enumerate}[leftmargin=1.5em,itemsep=0.1em]
\item \textbf{Mr.\ Pahlavi (President \& CEO)} --- final approval
on any executed SCLA, litigation filing, or strategic public
communication.
\item \textbf{Mr.\ Dupont (Managing Partner, Strategic Transactions;
largest investor)} --- deal lead. Conducts negotiation; selects
settlement tier; presents executed term sheet to Mr.\ Pahlavi.
\item \textbf{Mr.\ Donohue (General Counsel)} --- legal lead.
Drafts SCLA; supervises outside counsel; manages litigation hold;
leads pre-suit investigation; prepares SCLA for execution.
\item \textbf{Mr.\ Nandy (CTO; memorandum author)} --- operational and
technical support. Provides factual development on IP portfolio,
licensing record, and Counterparty-side operational context.
\end{enumerate}

\textbf{Recommended outside engagement:}

\begin{itemize}[leftmargin=1.5em,itemsep=0.1em]
\item \textbf{Outside litigation counsel}: a firm with DTSA + patent
+ Delaware-fiduciary-duty experience (Wilson Sonsini, Cooley, Fenwick
\& West, or Quinn Emanuel).
\item \textbf{Valuation firm}: independent fairness opinion on
equity-for-IP allocation (Houlihan Lokey, Duff \& Phelps / Kroll, or
specialist boutique).
\item \textbf{Investigative firm}: separate firm to investigate
Counterparty officer conduct (Kroll Investigations, FTI Consulting,
or specialist tech-IP investigators).
\end{itemize}

\textbf{Target close: 75 days} from Mr.\ Pahlavi's approval to
authorize the deal team, per the milestone schedule:

\begin{center}
\begin{tabular}{lp{10cm}}
\toprule
\textbf{Day} & \textbf{Action} \\
\midrule
0   & Mr.\ Pahlavi approves; outside counsel engaged. \\
7   & Litigation hold notice transmitted to Counterparty + named officers. \\
14  & Investigation firm begins document review. \\
21  & Term sheet delivered to Counterparty at selected tier. \\
30  & Counterparty disinterested-director response received. \\
45  & Investigation report finalized; settlement terms tightened/expanded. \\
60  & Definitive SCLA executed. \\
75  & Closing: cash payment, share issuance, license activation. \\
\bottomrule
\end{tabular}
\end{center}

% ─────────────────────────────────────────────────────────────────
\section*{XIV.\ Prayer for Relief (Template)}

In any litigation, the complaint's prayer will request:

\begin{enumerate}[leftmargin=1.5em,itemsep=0.1em]
\item Compensatory damages on each Count, anticipated \$80M to \$590M
on IP claims; \$180M to \$890M with opportunity-cost component.
\item Disgorgement of unjust enrichment from the current and any future raises (rebased to enterprise value, not round size).
\item Reasonable royalty (35 U.S.C.\ \S 284; 18 U.S.C.\ \S 1836(b)(3)(B)).
\item Exemplary damages 2$\times$ on DTSA (18 U.S.C.\ \S 1836(b)(3)(C)).
\item Treble damages on patent (35 U.S.C.\ \S 284).
\item Statutory damages on copyright (17 U.S.C.\ \S 504(c)) at elected
maximum.
\item Attorneys' fees (18 U.S.C.\ \S 1836(b)(3)(D); 17 U.S.C.\ \S 505;
35 U.S.C.\ \S 285; 15 U.S.C.\ \S 1117(a)).
\item Preliminary and permanent injunctive relief enjoining further
use absent executed commercial license.
\item Pre- and post-judgment interest.
\item Such other and further relief as the Court deems just.
\end{enumerate}

\textbf{Demand for jury trial} on all triable issues.

% ─────────────────────────────────────────────────────────────────
\vspace{1em}
\hrule height 0.4pt
\vspace{0.5em}
\begin{center}\textbf{\textsc{Schedules}}\end{center}
\vspace{0.5em}

\section*{Schedule A --- Lux IP Inventory in Production}

The full inventory is maintained at the companion license audit. Headline:

\begin{itemize}[leftmargin=1.5em,itemsep=0.1em]
\item BSD-3 / Apache-2.0 / MIT modules: approximately 25 (no Lux
action required).
\item LEL modules: \texttt{aml}, \texttt{broker}, \texttt{captable},
\texttt{compliance}, \texttt{corona}, \texttt{crypto}, \texttt{fhe},
\texttt{hsm}, \texttt{maker}, \texttt{mpc}, \texttt{precompile},
\texttt{ringtail}, \texttt{threshold}, \texttt{transfer},
\texttt{treasury}, \texttt{warp}.
\item LRL--PR modules: \texttt{vm}, \texttt{staking}, \texttt{sampler},
\texttt{dex} (Go bindings).
\item LGPL modules: \texttt{evm}.
\item Modules without LICENSE: \texttt{cex} and small number of
related; counsel should treat as no-rights-granted pending upstream
remediation.
\end{itemize}

\section*{Schedule B --- Lux IP Pipeline (Available Under SCLA)}

\begin{itemize}[leftmargin=1.5em,itemsep=0.1em]
\item AI inference and compute precompiles (LP-9080 family).
\item Advanced ZK rollup substrates.
\item Post-quantum hybrid signature schemes (X-Wing + ML-KEM; hybrid
ML-DSA + ECDSA).
\item Threshold MPC dynamic re-sharing protocols.
\item Confidential-execution attestation (TEE / SEV-SNP / TDX binding
to consensus).
\item White-label partner portal infrastructure.
\item Closed \texttt{\textasciitilde/work/luxcpp/*} (separate
Private-Moat Addendum required).
\end{itemize}

\section*{Schedule C --- Banking License Cost-to-Replicate}

\begin{center}
\begin{tabular}{lr}
\toprule
\textbf{License} & \textbf{Estimated Cost} \\
\midrule
US national or state banking charter & \$20--50M \\
Canadian Schedule II / partnership & \$10--20M \\
MAS Capital Markets Services + banking (Singapore) & \$5--10M \\
RBNZ deposit-taker (New Zealand) & \$5--15M \\
Sweden Finansinspektionen (EU passport) & \$5--10M \\
UK FCA EMI / Payment Institution & \$2--5M \\
Switzerland FINMA digital-asset banking & \$5--15M \\
Hong Kong SFC Type 1 / Type 7 + VATP & \$5--15M \\
\midrule
\textbf{Total direct regulatory + capital} & \textbf{\$57--140M} \\
Plus 5--7 years of timeline + execution risk & \\
\bottomrule
\end{tabular}
\end{center}

\section*{Schedule D --- Hanzo Inc.\ Portfolio Comparables}

\begin{itemize}[leftmargin=1.5em,itemsep=0.1em]
\item Damon Motorcycles --- electric / smart-motorcycle company.
\item Triller --- short-form video / social platform.
\item Bellabeat --- women's health wearables.
\item Unicorn Gold --- precious-metals tokenization platform.
\item Lux Network --- subject of this memorandum.
\item Zoo Labs Foundation / Zoo Network --- decentralized AI research network.
\end{itemize}

Retained-parent-equity range: 10\%--40\%. The 10\% ask sits at the
floor of that range.

\section*{Schedule E --- Mandatory Terms (Lux's Floor)}

See \S XI above. Sixteen mandatory terms, summarized:

\begin{enumerate}[leftmargin=1.5em,itemsep=0.05em]
\item 1 Lux director + 1 independent director seat (closing).
\item 99-year exclusive USA-market license.
\item Weighted-average broad-based anti-dilution + full ratchet
on down-rounds.
\item Pro-rata participation rights.
\item Tag-along and co-sale rights.
\item Information rights.
\item Patent grant (current and future Lux patents).
\item Maintenance and support of upstream OSS.
\item Reciprocal permanent license to Counterparty ATS / BD / TA /
listed securities / cross-listings.
\item Co-branding for Lux-network chains.
\item OnyxPlus ID verification at Lux internal cost.
\item AvaTrade routing into joint venue.
\item 8 bps crypto-conversion floor.
\item Warp precompile activation at next chain-upgrade.
\item Litigation hold and document preservation.
\item Independent investigation cooperation.
\end{enumerate}

\vfill
\hrule height 0.4pt
\vspace{0.5em}
\begin{small}
\textit{End of memorandum. This document is confidential, privileged,
and prepared in anticipation of litigation as well as for board
strategy purposes. The proposed settlement framework is non-binding
until executed as a definitive Settlement and Cross-License Agreement.
Distribution outside Lux Industries Inc., its Board of Directors, the
named decision-chain officers (Messrs.\ Pahlavi, Dupont, Donohue, and
Nandy), and engaged outside counsel requires the express written
authorization of Mr.\ Pahlavi and Mr.\ Donohue.}
\end{small}

\end{document}
